\subsection{特征标与格罗滕迪克群}

用 \(\theta_{G}\) 表示从 \(K\otimes_{\mathbf{Z}}\mathrm{R}_{\mathrm{K}}(\mathrm{G})\) 到 \(\mathcal{Z}_{\mathrm{K}}(\mathrm{G})\) 的 K-代数同态，它把 \(1\otimes[\pi]\) 映为 \(\chi_{\pi}\)（对每个有限维表示 \(\pi\)）。

命题 8. —— a) 同态 \(\theta_{\mathrm{G}}\) 是从 \(\mathrm{K} \otimes_{\mathbf{Z}} \mathrm{R}_{\mathrm{K}}(\mathrm{G})\) 到 \(\mathcal{Z}_{\mathrm{K}}(\mathrm{G})\) 的同构。

\begin{enumerate}
\def\labelenumi{\alph{enumi})}
\setcounter{enumi}{1}
\tightlist
\item
  假设 K 的特征为 0。那么 \(\theta_{\mathrm{G}}\) 定义了从 \(\mathrm{R}_{\mathrm{K}}(\mathrm{G})\) 到 \(\mathcal{Z}_{\mathrm{K}}(\mathrm{G})\) 的一个子环的同构，该子环由诸特征标 \(\chi_{\lambda}\)（\(\lambda\) 取遍 \(\widehat{\mathrm{G}}\)）的整系数线性组合组成。
\end{enumerate}

族 \(([\lambda])_{\lambda \in \widehat{\mathbf{G}}}\) 是 \(\mathbf{Z}\)-模 \(\mathrm{R}_{\mathrm{K}}(\mathrm{G})\) 的一组基，而族 \((\chi_{\lambda})_{\lambda \in \widehat{\mathbf{G}}}\) 是 K-向量空间 \(\mathcal{Z}_{\mathrm{K}}(\mathrm{G})\) 的一组基。断言 a) 与 b) 由此得出（VIII, 第 411 页，命题 5）。

\subsection{单表示的维数}

用 n 表示群 G 的基数。设 \(\pi\) 是 G 在有限维 K-向量空间 M 中的一个线性表示。对每个 \(g \in G\)，我们有 \(\pi(g)^{n} = 1_{M}\)，故 \(\pi(g)\) 的最小多项式整除 \(T^{n} - 1\)。由于在 K 中 \(n \cdot 1 \neq 0\)，这个最小多项式是可分的（V, §11, No.~2, 第 78 页）；又由于体 K 是代数闭的，M 的自同态 \(\pi(g)\) 可对角化（VII, §5, No.~7, 第 40 页，命题 12）。\(\pi(g)\) 的特征值都是 n 次单位根，并且对每个 \(\alpha \in K\)，\(\alpha\) 作为 \(\pi(g)\) 的特征值的几何重数（VII, §5, No.~2, 第 30 页，定义 1）等于 \(\alpha\) 作为 \(\pi(g)\) 的特征多项式的根的重数。用 \(\mathcal{O}_{n}\) 表示由 n 次单位根的集合 \(\mu_{n}(K)\) 生成的 K 的子群；它是一个有限生成的 Z-模，并且是 K 的一个子环。\(\pi\) 的特征标取值于 \(\mathcal{O}_{n}\)。

命题 9. —— 假设体 K 的特征为零。那么 G 的每个单表示的次数都整除 G 的基数。

设 \((\mathrm{V},\pi)\) 是 G 的一个单表示，\(\chi\) 是它的特征标。对 \(\mathrm{Z}(\mathrm{K}[\mathrm{G}])\) 的每个元素 a，V 的自同态 \(\pi(a)\) 是一个位似（VIII, 第 47 页，定理 1）；用 \(\varphi(a)\) 表示使得 \(\pi(a)=\varphi(a)_{\mathrm{V}}\) 的标量。

这样得到的映射 \(\varphi\)（从 \(Z(K[G])\) 到 \(K\) 的映射）是一个代数同态。取 \(a = \sum_{g \in G} \chi(g^{-1})g\)；由 VIII, 第 408 页的注，我们有 \(\varphi(a) = (\dim V)^{-1}|G|\)。另一方面，\(a\) 属于 \(\mathcal{O}_n[G] \cap Z(K[G])\)，即 \(K[G]\) 的一个子环，后者是一个有限生成的 \(\mathbf{Z}\)-模（VII, §3, 第 15 页，推论）。因此元素 \(\varphi(a) = (\dim V)^{-1}|G|\)（\(K\) 的元素）属于 \(K\) 的一个作为有限生成 \(\mathbf{Z}\)-模的子环。我们借助下面的引理来完成证明。

引理. —— 设 L 是 Q 的一个扩张，A 是 L 的一个子环。假设 A 是有限生成的 Z-模。我们有 \(A \cap Q = Z\)。

由于 Z-模 \(A \cap Q\) 是有限生成的，存在严格正整数 N，使得 \(A \cap Q\) 包含于 \(\frac{1}{N}Z\)。设 x 是 \(Q - Z\) 的一个元素；把它写成 \(x = \frac{p}{q}\)，其中 p 与 q 是互素的整数且 \(q \geqslant 2\)。我们有 \(q^{N} \geqslant 2^{N} > N\)（集合论 III, §3, No.~6, 第 165 页，定理 2），整数 \(p^{N}\) 与 \(q^{N}\) 互素，从而 \(x^{N} \notin \frac{1}{N}Z\)。由此可知 x 不属于 A。引理证毕。

在下一小节中，我们把命题 9 推广到只假设 K 的特征不整除 G 的阶的情形。

\subsection{换基域}

我们沿用上一小节的记号。设 \(K'\) 是一个代数闭体，使得元素 \(n \cdot 1\) 在 \(K'\) 中非零。群 \(\mu_{n}(\mathrm{K})\) 与 \(\mu_{n}(\mathrm{K}')\) 都是 n 阶循环群（V, §11, No.~2, 第 78 页，定理 1）。取定一个同构 \(\varphi\)（从 \(\mu_{n}(\mathrm{K})\) 到 \(\mu_{n}(\mathrm{K}')\) 的同构）。设 \(\pi\) 是 G 在有限维 K-向量空间中的一个线性表示，并设 \(\pi'\) 是 G 在有限维 \(K'\)-向量空间中的一个线性表示。我们称 \(\pi\) 与 \(\pi'\)（通过 \(\varphi\)）相关联，如果对每个 \(g \in G\) 与每个 \(\omega \in \mu_{n}(\mathrm{K})\)，\(\omega\) 作为 \(\pi(g)\) 的特征值的重数等于 \(\varphi(\omega)\) 作为 \(\pi'(g)\) 的特征值的重数。这时 \(\pi\) 与 \(\pi'\) 有相同的维数，取 g = 1 即可看出。

设 \(\pi_1\) 与 \(\pi_2\)（相应地，\(\pi_1'\) 与 \(\pi_2'\)）是 \(G\) 在 \(K\)（相应地，\(K'\)）上有限维向量空间中的线性表示。我们有下列性质：

\begin{enumerate}
\def\labelenumi{\alph{enumi})}
\item
  若 \(\pi_1\) 与 \(\pi_1'\) 和 \(\pi_2'\) 都相关联，则 \(\pi_1'\) 与 \(\pi_2'\) 同构。
\item
  若 \(\pi_1\) 与 \(\pi_1'\) 相关联且 \(\pi_2\) 与 \(\pi_2'\) 相关联，则 \(\pi_1 \oplus \pi_2\) 与 \(\pi_1' \oplus \pi_2'\) 相关联，且 \(\pi_1 \otimes \pi_2\) 与 \(\pi_1' \otimes \pi_2'\) 相关联。
\end{enumerate}

断言 a) 由 VIII, 第 402 页的推论 5 得出，而断言 b) 是显然的。

这里，\(\mathcal{S}_{\mathrm{K}}(\mathrm{G})\) 表示单 K{[}G{]}-模的类所成的集合（先前记作 \(\widehat{G}\)）；同样定义 \(\mathcal{S}_{\mathrm{K}^{\prime}}(\mathrm{G})\)。集合 \(\mathcal{S}_{\mathrm{K}}(\mathrm{G})\) 与 \(\mathcal{S}_{\mathrm{K}^{\prime}}(\mathrm{G})\) 都是有限的，其基数为共轭类的个数（VIII, 第 411 页，命题 5）。

命题 10. —— 存在唯一的映射 \(\varphi_{\mathrm{G}}\)，从 \(\mathcal{S}_{\mathrm{K}}(\mathrm{G})\) 到 \(\mathcal{S}_{\mathrm{K}'}(\mathrm{G})\)，使得 \(\lambda\) 与 \(\varphi_{\mathrm{G}}(\lambda)\) 通过 \(\varphi\) 相关联——这里 \(\lambda\) 取遍 \(\mathcal{S}_{\mathrm{K}}(\mathrm{G})\)。此外，\(\varphi_{\mathrm{G}}\) 是双射。

\(\varphi_{G}\) 的唯一性由上面的性质 a) 得出。

\textbf{A) 假设体 K 的特征为 0。}

群 \(\mu_n(\mathrm{K})\) 是循环群（V, §11, No.~2, 第 78 页，定理 1）；取该群的一个生成元 \(\zeta\)。考虑环同态 \(\rho: \mathbf{Z}[\mathrm{X}] \to \mathcal{O}_n\)，它把 X 映为 \(\zeta\)。它是满射。分圆多项式 \(\Phi_n(\mathrm{X})\) 在 \(\mathbf{Q}[\mathrm{X}]\) 中不可约（V, §11, No.~5, 第 84 页，定理 2）；因此它是 \(\zeta\) 在 \(\mathbf{Q}\) 上的最小多项式。多项式 \(\Phi_n\) 是首一的且具有整系数（V, §11, No.~4, 第 81 页）。设 \(\mathrm{P} \in \mathbf{Z}[\mathrm{X}]\) 是满足 \(\mathrm{P}(\zeta) = 0\) 的一个多项式；由多项式的欧几里得除法（IV, §1, No.~6, 第 10 页），存在 \(\mathbf{Z}[\mathrm{X}]\) 中的两个多项式 Q 与 R，使得 \(\mathrm{P} = \mathrm{Q}\Phi_n + \mathrm{R}\) 且 \(\deg(\mathrm{R}) < \deg(\Phi_n)\)。我们有 \(\mathrm{R}(\zeta) = 0\)，从而 \(\mathrm{R} = 0\)，因为 \(\Phi_n\) 是 \(\zeta\) 的最小多项式。于是 \(\rho\) 的核是理想 \(\Phi_n\mathbf{Z}[\mathrm{X}]\)（\(\mathbf{Z}[\mathrm{X}]\) 中的理想），而 \(\rho\) 诱导出从 \(\mathbf{Z}[\mathrm{X}] / \Phi_n\mathbf{Z}[\mathrm{X}]\) 到 \(\mathcal{O}_n\) 的环同构。

令 \(\zeta' = \varphi(\zeta)\)；它是一个 \(n\) 次本原单位根，位于 \(\mathbf{K}'\) 中，因此我们有 \(\Phi_n(\zeta') = 0\)（V, §11, No.~5, 第 83 页，引理 3）。从而存在同态 \(\varphi_0\)（从环 \(\mathcal{O}_n\) 到体 \(\mathbf{K}'\) 的同态），它把 \(\zeta\) 变为 \(\zeta'\)；它是映射 \(\varphi\)（从 \(\mu_n(\mathrm{K})\) 到 \(\mu_n(\mathrm{K}')\) 的映射）的扩张。设 \(\mathcal{O}\) 是 \(\mathrm{K}\) 中由形如 \(\frac{a}{n^r}\)（\(a \in \mathcal{O}_n\)，\(r \in \mathbf{N}\)）的元素组成的子环。由于 \(n \cdot 1\) 在 \(\mathbf{K}'\) 中可逆，同态 \(\varphi_0\) 可扩张为一个同态 \(\varphi_1\)（从 \(\mathcal{O}\) 到 \(\mathbf{K}'\) 的同态）。

我们把 \(\mathcal{O}[\mathrm{G}]\)（群 \(\mathrm{G}\) 在 \(\mathcal{O}\) 上的代数）与代数 \(\mathrm{K}[\mathrm{G}]\) 的一个子环等同起来，并借助下述公式定义环同态 \(\Phi\)（从 \(\mathcal{O}[\mathrm{G}]\) 到 \(\mathrm{K}'[\mathrm{G}]\) 的环同态）：

\[
\Phi \bigg (\sum_ {g \in \mathrm{G}} a _ {g} g \bigg) = \sum_ {g \in \mathrm{G}} \varphi_ {1} (a _ {g}) g.\tag{42}
\]

用 \(\mathcal{C}\) 表示 G 的共轭类所成的集合。对 \(\mathcal{C}\) 中的 C，用 \(u_{\mathrm{C}}\) 表示元素 \(\sum_{g\in \mathrm{C}}g\)（\(\mathcal{O}[\mathrm{G}]\) 的元素）；族 \((u_{\mathrm{C}})_{\mathrm{C}\in \mathcal{C}}\) 构成 \(\mathcal{O}\) 上向量空间 \(\mathrm{Z}(\mathcal{O}[\mathrm{G}])\)（代数 \(\mathcal{O}[\mathrm{G}]\) 的中心）的一组基（VIII, 第 398 页）。对任意元素 \(\lambda\)（\(\mathcal{S}_{\mathrm{K}}(\mathrm{G})\) 的元素），用 \(\chi_{\lambda}\) 表示它的特征标，用 \(d_{\lambda}\) 表示它的维数，并用 \(e_{\lambda}\) 表示

由下式定义的 K{[}G{]} 中的元素

\[
e _ {\lambda} = | \mathrm{G} | ^ {- 1} d _ {\lambda} \sum_ {g \in \mathrm{G}} \chi_ {\lambda} (g ^ {- 1}) g.\tag{43}
\]

族 \((e_{\lambda})\) 是 \(Z(K[G])\)（环 K{[}G{]} 的中心）在 K 上的一组基（VIII, 第 408 页）。我们有

\[
e _ {\lambda} = \sum_ {\mathrm{C} \in \mathcal {C}} \alpha_ {\lambda , \mathrm{C}} u _ {\mathrm{C}}\tag{44}
\]

其中

\[
\alpha_ {\lambda , \mathrm{C}} = | \mathrm{G} | ^ {- 1} d _ {\lambda} \chi_ {\lambda} (\mathrm{C} ^ {- 1}).\tag{45}
\]

对 \(C \in \mathcal{C}\) 与 \(\lambda \in \mathcal{S}_{\mathrm{K}}(\mathrm{G})\)，令 \(\beta_{\mathrm{C},\lambda} = |\mathrm{G}| d_{\lambda}^{-1} d(\mathrm{C})^{-1} \chi_{\lambda}(\mathrm{C})\)。矩阵 \((\alpha_{\lambda,\mathrm{C}})\) 的元素都属于 \(\mathcal{O}\)，并且由 VIII, 第 411 页的公式 (31)，它的逆矩阵是矩阵 \((\beta_{\mathrm{C},\lambda})\)，而后者的元素由 VIII, 第 415 页的命题 9 也都属于 \(\mathcal{O}\)。于是族 \((e_{\lambda})\) 是 \(\mathcal{O}\)-模 \(\mathrm{Z}(\mathcal{O}[\mathrm{G}])\) 的一组基。

诸元素 \(\Phi(u_{\mathrm{C}}) = \sum_{g\in \mathrm{C}}g\)（\(K^{\prime}[G]\) 中的元素）构成 \(K^{\prime}\) 上向量空间 \(Z(K^{\prime}[G])\)（环 \(K^{\prime}[G]\) 的中心）的一组基。我们有

\[
\Phi (e _ {\lambda}) = \sum_ {\mathrm{C} \in \mathcal {C}} \varphi_ {1} (\alpha_ {\lambda , \mathrm{C}}) \Phi (u _ {\mathrm{C}}),\tag{46}
\]

并且以 \(\varphi_{1}(\alpha_{\lambda,\mathrm{C}})\) 为元素的矩阵可逆。因此诸 \(\Phi(e_{\lambda})\) 所成的族是 \(\mathrm{Z}(\mathrm{K}^{\prime}[\mathrm{G}])\) 的一组基。族 \((e_{\lambda})\) 是 \(\mathrm{Z}(\mathrm{K}[\mathrm{G}])\) 中幂等元 1 的一个分拆（VIII, 第 146 页与第 408 页）；换言之，我们有

\[
\sum_ {\lambda} e _ {\lambda} = 1, \quad e _ {\lambda} ^ {2} = e _ {\lambda}, \quad e _ {\lambda} e _ {\mu} = 0 \quad \mathrm{if} \quad \lambda \neq \mu .
\]

由此推出，诸 \(\Phi(e_{\lambda})\) 所成的族是 \(Z(K^{\prime}[G])\) 中幂等元 1 的一个分拆；由于该族是 \(Z(K^{\prime}[G])\) 在 \(K^{\prime}\) 上的一组基，其元素是 \(Z(K^{\prime}[G])\) 中的不可分解幂等元（VIII, 第 148 页，注 4）。

对 \(\lambda'\) 属于 \(\mathcal{S}_{\mathrm{K}'}(\mathrm{G})\) 的情形，如上定义 \(\chi_{\lambda'}, d_{\lambda'}\) 与 \(e_{\lambda'}\)。由 VIII, 第 408 页，诸元素 \(e_{\lambda'}\) 是 \(\mathrm{Z}(\mathrm{K}'[\mathrm{G}])\) 中的不可分解幂等元。因此存在双射 \(\varphi_{\mathrm{G}}\)，从 \(\mathcal{S}_{\mathrm{K}}(\mathrm{G})\) 到 \(\mathcal{S}_{\mathrm{K}'}(\mathrm{G})\)，使得 \(\Phi(e_{\lambda}) = e_{\varphi_{\mathrm{G}}(\lambda)}\)，其中 \(\lambda\) 取遍 \(\mathcal{S}_{\mathrm{K}}(\mathrm{G})\)。

设 \(\lambda\) 是 \(\mathcal{S}_{\mathrm{K}}(\mathrm{G})\) 中的元素；令 \(\lambda' = \varphi_{\mathrm{G}}(\lambda)\)。设 \((\mathrm{V}_{\lambda}, \pi_{\lambda})\)（相应地，\((\mathrm{V}_{\lambda'}, \pi_{\lambda'})\)）是 G 的一个线性表示，其相伴的 K{[}G{]}-模（相应地，K'{[}G{]}-模）的类为 \(\lambda\)（相应地，\(\lambda'\)）。我们来证明 \(\lambda\) 与 \(\lambda'\) 相关联。设 \(g\) 是 G 的一个元素。设 \(\delta(\mathrm{T})\) 是自同态 \(1 + \mathrm{T}\pi_{\lambda}(g)\)

（K{[}T{]}-模 K{[}T{]} \(\otimes_{\mathrm{K}}\mathrm{V}_{\lambda}\) 的自同态）的行列式。用 \(\omega_{1},\ldots ,\omega_{d_{\lambda}}\) 表示 \(\pi_{\lambda}(g)\) 的特征值；我们有

\[
\delta (\mathrm{T}) = (1 + \mathrm{T} \omega_ {1}) \dots (1 + \mathrm{T} \omega_ {d _ {\lambda}}).
\]

我们同样定义 \(\delta'(T)\)，并把 \(\pi_{\lambda'}(g)\) 的特征值记作 \(\omega_1', \ldots, \omega_{d_{\lambda'}}'\)。\(\mathcal{O}\)-模 \(\mathcal{O}[G]\) 是自由的，以 \(G\) 为基，而 \(K\)-向量空间 \(K[G]\) 以 \(G\) 为基。用 \(\Delta(T)\) 表示乘以 \(1 + e_{\lambda}gT\)（在 \(\mathcal{O}[T]\)-模 \(\mathcal{O}[T] \otimes_{\mathcal{O}} \mathcal{O}[G]\) 中）的行列式。它也是乘以 \(1 + e_{\lambda}gT\)（在 \(K[T]\)-模 \(K[T] \otimes_K K[G]\) 中）的行列式。设 \(\overline{\varphi}_1\) 是从 \(\mathcal{O}[T]\) 到 \(K'[T]\) 的同态，它扩张 \(\varphi_1\) 且把 \(T\) 映为 \(T\)。由于 \(G\) 是 \(K'\)-向量空间 \(K'[G]\) 的一组基，多项式 \(\overline{\varphi}_1(\Delta(T))\) 等于行列式 \(\Delta'(T)\)，即乘以 \(1 + e_{\lambda'}gT\)（在 \(K'\)-向量空间 \(K'[G]\) 中）的行列式。

代数 K{[}G{]} 是它的诸单分量 \(e_{\mu}K[G]\)（\(\mu\) 取遍 \(\mathcal{S}_{\mathrm{K}}(\mathrm{G})\)）的直和。当 \(\mu\) 异于 \(\lambda\) 时，元素 \(e_{\lambda}g\) 零化 \(e_{\mu}K[G]\)。此外，乘以 \(e_{\lambda}g\) 的乘法与乘以 g 的乘法在 \(e_{\lambda}K[G]\) 中一致。注意到 VIII, 第 409 页以及 VIII, 第 400 页的例 6，G 在 \(e_{\lambda}K[G]\) 中的表示是 \(d_{\lambda}\) 个类为 \(\lambda\) 的表示的直和。因此我们有 \(\Delta(\mathrm{T}) = \delta(\mathrm{T})^{d_{\lambda}}\)。

同理，我们有 \(\Delta^{\prime}(\mathrm{T}) = \delta^{\prime}(\mathrm{T})^{d_{\lambda^{\prime}}}\)。

由关系式 \(\Delta'(\mathrm{T}) = \overline{\varphi}_1(\Delta(\mathrm{T}))\)，我们首先推出 \(d_{\lambda}^2 = d_{\lambda'}^2\)，从而 \(d_{\lambda} = d_{\lambda'}\)；然后推出序列 \(\varphi(\omega_1), \ldots, \varphi(\omega_{d_\lambda})\) 可由序列 \((\omega_1', \ldots, \omega_{d_{\lambda'}'}')\) 经过指标集的一个置换得到。

由于这一点对每个元素 \(g\)（\(G\) 的元素）都成立，表示 \(\lambda\) 与 \(\lambda'\) 相关联。这样，当体 \(K\) 的特征为 0 时，我们就证明了命题 10。

\textbf{B) 一般情形。}

设 L 是一个特征为 0 的代数闭体（例如 Q 的一个代数闭包）。用 \(\mathcal{S}_{\mathrm{L}}(\mathrm{G})\) 表示单 L{[}G{]}-模的类所成的集合。取定一个同构 \(\eta\)（从群 \(\mu_{n}(\mathrm{L})\) 到群 \(\mu_{n}(\mathrm{K})\) 的同构），并令 \(\eta' = \varphi \circ \eta\)。由证明的部分 A)，存在双射

\[
\eta_ {\mathrm{G}}: \mathcal {S} _ {\mathrm{L}} (\mathrm{G}) \to \mathcal {S} _ {\mathrm{K}} (\mathrm{G}), \quad \eta_ {\mathrm{G}} ^ {\prime}: \mathcal {S} _ {\mathrm{L}} (\mathrm{G}) \to \mathcal {S} _ {\mathrm {K^ {\prime}}} (\mathrm{G})
\]

它们具有下述性质：对每个 \(\lambda\)（\(\mathcal{S}_{\mathrm{L}}(\mathrm{G})\) 中的元素），表示 \(\lambda\) 与 \(\eta_{\mathrm{G}}(\lambda)\) 通过 \(\eta\) 相关联，且表示 \(\lambda\) 与 \(\eta_{\mathrm{G}}^{\prime}(\lambda)\) 通过 \(\eta^{\prime}\) 相关联。双射 \(\varphi_{G} = \eta_{G}^{\prime} \circ \eta_{G}^{-1}\) 具有所要的性质。

双射 \(\varphi_{G}\)（从 \(\mathcal{S}_{K}(G)\) 到 \(\mathcal{S}_{K'}(G)\) 的双射）可扩张为一个同构（仍记作 \(\varphi_{G}\)），即从格罗滕迪克群 \(R_{K}(G)\) 到群 \(R_{K'}(G)\) 的同构。

注 1. —— 假设 \(K'\) 是 K 的一个扩张，且同构 \(\varphi\) 是映射 \(\xi \mapsto \xi \cdot 1\)；那么映射 \(\varphi_{G}\) 由从 K 到 \(K'\) 的标量扩张给出。

推论 1. —— 映射 \(\varphi_{\mathrm{G}}\) 是从 \(\mathrm{R}_{\mathrm{K}}(\mathrm{G})\) 到 \(\mathrm{R}_{\mathrm{K}^{\prime}}(\mathrm{G})\) 的环同构。对每个有限维表示 \(\pi\)（\(\mathrm{G}\) 在 K-向量空间中的表示），我们有 \(\varphi_{\mathrm{G}}([\pi]) = [\pi^{\prime}]\)，其中 \(\pi^{\prime}\) 是与 \(\pi\) 通过 \(\varphi\) 相关联的一个表示。

这由 G 的表示的半单性以及 VIII, 第 416 页的性质 b) 得出。

推论 2. —— G 的每个单表示的维数整除 G 的阶。

这由命题 10 以及 VIII, 第 415 页的命题 9 得出。

注. —— 2) 假设群 G 是交换群。我们在 VIII, 第 414 页的注中已看到，\(\mathcal{S}_{\mathrm{K}}(\mathrm{G})\) 可与集合 \(\operatorname{Hom}(\mathrm{G}, \mu_n(\mathrm{K}))\) 等同。同样，\(\mathcal{S}_{\mathrm{K}'}(\mathrm{G})\) 可与 \(\operatorname{Hom}(\mathrm{G}, \mu_n(\mathrm{K}'))\) 等同。在这些等同之下，双射 \(\varphi_{\mathrm{G}}\) 就是映射 \(\chi \mapsto \varphi \circ \chi\)。

\begin{enumerate}
\def\labelenumi{\arabic{enumi})}
\setcounter{enumi}{2}
\tightlist
\item
  设 \(\pi_1\) 与 \(\pi_2\) 是 G 在 K 上有限维向量空间中的线性表示。对 \(i = 1,2\)，设 \(\pi_i'\) 是与 \(\pi_i\) 通过 \(\varphi\) 相关联的一个表示。我们有
\end{enumerate}

\[
\dim_ {K} \operatorname{Hom} _ {K} (\pi_ {1}, \pi_ {2}) = \dim_ {K ^ {\prime}} \operatorname{Hom} _ {K ^ {\prime}} (\pi_ {1} ^ {\prime}, \pi_ {2} ^ {\prime}).\tag{47}
\]

证明沿用 VIII, 第 410 页的推论的证明，归结为诸 \(\pi_{i}\)（从而诸 \(\pi_{i}^{\prime}\)）都是单表示的情形。

\begin{enumerate}
\def\labelenumi{\arabic{enumi})}
\setcounter{enumi}{3}
\tightlist
\item
  设 H 是 G 的一个基数为 m 的子群。同构 \(\varphi\) 限制为从 \(\mu_{m}(\mathrm{K})\) 到 \(\mu_{m}(\mathrm{K}^{\prime})\) 的同构，从而给出一个环同构 \(\varphi_{H}\)，从 \(\mathrm{R}_{\mathrm{K}}(\mathrm{H})\) 到 \(\mathrm{R}_{\mathrm{K}^{\prime}}(\mathrm{H})\)。下面的图交换：
\end{enumerate}

\[
\begin{array}{c c} \mathrm {R_ {K} (G)} \xrightarrow {\mathrm {Res_ {H} ^ {G}}} \mathrm {R_ {K} (H)} & \qquad \qquad \mathrm {R_ {K} (H)} \xrightarrow {\mathrm {Ind_ {H} ^ {G}}} \mathrm {R_ {K} (G)} \\ \Biggl \downarrow \varphi_ {\mathrm{G}} & \Biggl \downarrow \varphi_ {\mathrm{H}} \\ \mathrm {R_ {K^ {\prime}} (G)} \xrightarrow {\mathrm {Res_ {H} ^ {G}}} \mathrm {R_ {K^ {\prime}} (H)}, & \qquad \qquad \mathrm {R_ {K^ {\prime}} (H)} \xrightarrow {\mathrm {Ind_ {H} ^ {G}}} \mathrm {R_ {K^ {\prime}} (G)}. \end{array}
\]

第一个图的交换性是显然的，而第二个图的交换性由它借助弗罗贝尼乌斯互反律与公式 (47) 得出。

\begin{enumerate}
\def\labelenumi{\arabic{enumi})}
\setcounter{enumi}{4}
\tightlist
\item
  假设 \(G\) 是两个有限群的乘积 \(G' \times G''\)。我们按上例的方式定义同构 \(\varphi_{G'}\) 与 \(\varphi_{G''}\)。
\end{enumerate}

我们有一个交换图

\[
\begin{array}{c} \mathrm{R} _ {\mathrm{K}} (\mathrm{G} ^ {\prime}) \otimes_ {\mathbf {Z}} \mathrm{R} _ {\mathrm{K}} (\mathrm{G} ^ {\prime \prime}) \xrightarrow {\kappa} \mathrm{R} _ {\mathrm{K}} (\mathrm{G}) \\ \Biggl \downarrow \varphi_ {\mathrm{G} ^ {\prime}} \otimes \varphi_ {\mathrm{G} ^ {\prime \prime}} \qquad \qquad \qquad \Biggl \downarrow \varphi_ {\mathrm{G}} \\ \mathrm{R} _ {\mathrm{K} ^ {\prime}} (\mathrm{G} ^ {\prime}) \otimes_ {\mathbf {Z}} \mathrm{R} _ {\mathrm{K} ^ {\prime}} (\mathrm{G} ^ {\prime \prime}) \xrightarrow {\kappa^ {\prime}} \mathrm{R} _ {\mathrm{K} ^ {\prime}} (\mathrm{G})  , \end{array}
\]

其中同构 \(\kappa\) 与 \(\kappa'\) 是 VIII, 第 406 页中定义的那两个同构。

\subsection{复线性表示}

在本小节中，我们假设 K 是复数域 C。

设 \((M,\pi)\) 是 G 的一个线性表示。若 M 上的埃尔米特型 \(\Phi\) 满足

\[
\Phi (\pi (g) x, \pi (g) x ^ {\prime}) = \Phi (x, x ^ {\prime})\tag{48}
\]

（对所有 \(x, x' \in M\) 与每个 \(g \in G\)），则称它在 G 下不变。这也意味着，对每个 \(g \in G\)，M 的自同构 \(\pi(g)\) 关于 \(\Phi\) 是酉的。

命题 11. —— 设 \((M,\pi)\) 是 G 的一个有限维线性表示。

\begin{enumerate}
\def\labelenumi{\alph{enumi})}
\item
  在 M 上存在一个正的、分离的且在 G 下不变的埃尔米特型。
\item
  假设表示 \(\pi\) 是单表示。若 \(\Phi\) 与 \(\Psi\) 是 M 上非零的且在 G 下不变的埃尔米特型，则存在实数 a 使得 \(\Psi = a\Phi\)。
\end{enumerate}

在向量空间 M 上选取一个分离的正埃尔米特型，并把它记作 \(\Phi_{0}\)。我们通过下面的式子定义一个在 G 下不变的分离的正埃尔米特型 \(\Phi\)：

\[
\Phi (x, x ^ {\prime}) = \sum_ {g \in \mathrm{G}} \Phi_ {0} (\pi (g) x, \pi (g) x ^ {\prime})\tag{49}
\]

（对 \(x, x' \in M\)）。

设 \(\Psi\) 是 M 上的一个埃尔米特型；存在 M 的唯一的自同态 A，使得 \(\Psi(x, x') = \Phi(x, Ax')\)（对 M 中的 \(x, x'\)）。此外，若 \(\Psi\) 在 G 下不变，则自同态 A 与自同构 \(\pi(g)\)（对 \(g \in G\)）交换。若表示 \(\pi\) 是单表示，则由舒尔引理（VIII, 第 47 页，定理 1），A 是一个位似，从而存在复数 \(a\) 使得 \(\Psi = a\Phi\)。由于 \(\Phi\) 与 \(\Psi\) 都是埃尔米特型且 \(\Phi\) 非零，\(a\) 是实数。命题得证。

我们赋予向量空间 \(\mathbf{C}[\mathrm{G}]\)（\(\mathrm{G}\) 上复函数所成的向量空间）一个希尔伯特空间结构，其内积由

\[
\langle f | f ^ {\prime} \rangle_ {\mathrm{G}} = | \mathrm{G} | ^ {- 1} \sum_ {g \in \mathrm{G}} \overline {{f (g)}} f ^ {\prime} (g).\tag{50}
\]

给出。对任意函数 \(f \in C[G]\)，用 \(f^{*}\) 表示由

\[
f ^ {*} (g) = \overline {{f (g ^ {- 1})}}\tag{51}
\]

（对 \(g \in \mathbf{G}\)）定义的函数；映射 \(f \mapsto f^*\) 是 \(\mathbf{C}[\mathbf{G}]\) 的一个半线性对合。我们还有

\[
\langle f | f ^ {\prime} \rangle_ {\mathrm{G}} = \langle f ^ {*}, f ^ {\prime} \rangle_ {\mathrm{G}}\tag{52}
\]

（对 \(f, f' \in \mathbf{C}[\mathrm{G}]\)），这里采用 VIII, 第 410 页公式 (28) 的记号。因此我们有

\[
\langle f | f ^ {\prime} \rangle_ {\mathrm{G}} = | \mathrm{G} | ^ {- 2} \tau (f ^ {*} f ^ {\prime}).\tag{53}
\]

设 \((\mathrm{M},\pi)\) 是 G 的一个有限维线性表示。我们赋予向量空间 M 一个使诸自同态 \(\pi(g)\) 为酉自同构的希尔伯特空间结构（命题 11）。若用 \(A^{*}\) 表示 M 的自同态 A 关于这个结构的伴随，则我们有 \(\mathrm{Tr}(A^{*})=\overline{\mathrm{Tr}(A)}\)。对每个 \(g\in G\)，我们有 \(\pi(g^{-1})=\pi(g)^{*}\)，从而 \(\chi_{\pi}(g^{-1})=\overline{\chi_{\pi}(g)}\)；换言之，我们有 \(\chi_{\pi}=\chi_{\pi}^{*}\)。于是特征标的正交关系（VIII, 第 410 页，命题 4）取如下形式

\[
\langle \chi_ {\lambda} | \chi_ {\mu} \rangle_ {\mathrm{G}} = \delta_ {\lambda \mu}\tag{54}
\]

（对 \(\lambda, \mu \in \widehat{\mathbf{G}}\)）。它表达了这样一个事实：特征标族 \((\chi_{\lambda})_{\lambda \in \widehat{\mathbf{G}}}\)——\(\mathbf{G}\) 的单表示的特征标族——是由中心函数组成的希尔伯特空间 \(\mathrm{Z}(\mathbf{C}[\mathrm{G}])\) 的一组规范正交基。

设 \(\pi\) 与 \(\pi'\) 是 G 的有限维线性表示。我们有关系式 \(\langle \chi_{\pi}|\chi_{\pi'}\rangle_{\mathrm{G}} = \dim_{\mathbf{C}}\mathrm{Hom}_{\mathrm{G}}(\pi, \pi')\)（VIII, 第 410 页，推论）。表示 \(\pi\) 不可约，当且仅当 \(\langle \chi_{\pi}|\chi_{\pi}\rangle_{\mathrm{G}} = 1\)。

对每个元素 \(\lambda\)（\(\widehat{G}\) 的元素），我们赋予向量空间 \(V_{\lambda}\) 一个使诸自同构 \(\pi_{\lambda}(g)\) 为酉自同构的希尔伯特空间结构。用 \(\langle v|v'\rangle_{\lambda}\) 表示两个元素 \(v, v'\)（\(V_{\lambda}\) 的元素）的内积，用 \(u^{*}\) 表示自同态 \(u\)（\(V_{\lambda}\) 的自同态）的伴随。设 \(A = (A_{\lambda})_{\lambda \in \widehat{G}}\) 与 \(A' = (A_{\lambda}')_{\lambda \in \widehat{G}}\) 是 \(\mathrm{F}(\widehat{\mathrm{G}})\) 的元素。我们记 \(\mathrm{A}^* = (\mathrm{A}_\lambda^*)_{\lambda \in \widehat{\mathrm{G}}}\)。我们有 \(\overline{\mathcal{F}} (a^{*}) = (\overline{\mathcal{F}} (a))^*\)（对每个元素 \(a\)（\(\mathbf{C}[\mathrm{G}]\) 中的元素））。令

\[
\langle \mathrm{A} | \mathrm{A} ^ {\prime} \rangle_ {\widehat {\mathrm{G}}} = | \mathrm{G} | ^ {- 2} \widehat {\tau} (\mathrm{A} ^ {*} \mathrm{A} ^ {\prime}).\tag{55}
\]

由 VIII, 第 407 页的公式 (10)，我们有

\[
\langle \mathrm{A} | \mathrm{A} ^ {\prime} \rangle_ {\widehat {\mathrm{G}}} = \frac {1}{| \mathrm{G} | ^ {2}} \sum_ {\lambda \in \widehat {\mathrm{G}}} d _ {\lambda} \operatorname{Tr} (\mathrm{A} _ {\lambda} ^ {*} \mathrm{A} _ {\lambda} ^ {\prime}).\tag{56}
\]

由于 \(\hat{\tau}\circ\overline{{\mathcal{F}}}=\tau\)，公式 (53) 与 (55) 蕴含映射 \(\overline{{\mathcal{F}}}\) 是从 C{[}G{]} 到 F( \(\widehat{G}\) ) 的希尔伯特空间同构。

舒尔正交关系（VIII, 第 410 页）可以用希尔伯特内积重新表述。于是关系式 (21) 与 (24) 给出下述断言。对 \(\lambda \in \widehat{\mathbf{G}}\) 与 \(x, x', y, y'\)（\(\mathrm{V}_{\lambda}\) 中的元素），我们有

\[
| \mathrm{G} | ^ {- 1} \sum_ {g \in \mathrm{G}} \overline {{\langle x | \pi_ {\lambda} (g) x ^ {\prime} \rangle}} _ {\lambda} \langle y | \pi_ {\lambda} (g) y ^ {\prime} \rangle_ {\lambda} = d _ {\lambda} ^ {- 1} \overline {{\langle x | y \rangle}} _ {\lambda} \langle x ^ {\prime} | y ^ {\prime} \rangle_ {\lambda}.\tag{57}
\]

若 \(\lambda\) 与 \(\mu\) 是 \(\widehat{\mathrm{G}}\) 的两个不同元素，则对 \(x, x'\)（\(\mathrm{V}_{\lambda}\) 中的元素）与 \(y, y'\)（\(\mathrm{V}_{\mu}\) 中的元素），我们有

\[
\sum_ {g \in \mathrm{G}} \overline {{\langle x | \pi_ {\lambda} (g) x ^ {\prime} \rangle_ {\lambda}}} \langle y | \pi_ {\mu} (g) y ^ {\prime} \rangle_ {\mu} = 0.\tag{58}
\]

对每个 \(\lambda \in \widehat{\mathbf{G}}\)，取定规范正交基 \((e_{\lambda,i})_{1\leqslant i\leqslant d_{\lambda}}\)（\(\mathrm{V}_{\lambda}\) 的一组规范正交基）。对任意 \(g\in \mathbf{G}\)，用 \((\pi_{ij}^{\lambda}(g))\) 表示自同态 \(\pi_{\lambda}(g)\)（\(\mathrm{V}_{\lambda}\) 的自同态）关于这组基的矩阵；我们有

\[
\pi_ {i j} ^ {\lambda} (g) = \langle e _ {\lambda , i} | \pi_ {\lambda} (g) e _ {\lambda , j} \rangle_ {\lambda}.\tag{59}
\]

由于自同态 \(\pi_{\lambda}(g)\) 是酉的，它的逆等于 \(\pi_{\lambda}(g)^*\)，于是

\[
\overline {{\pi_ {i j} ^ {\lambda} (g)}} = \pi_ {j i} ^ {\lambda} (g ^ {- 1}).\tag{60}
\]

于是由 VIII, 第 409 页的公式 (22) 以及第 409 页的 (25) 可知，诸函数 \((d_{\lambda})^{1 / 2}\pi_{ij}^{\lambda}\)（\(\lambda \in \widehat{\mathbf{G}}\)，\(1\leqslant i\leqslant d_{\lambda}\)，\(1\leqslant j\leqslant d_{\lambda}\)）构成希尔伯特空间 \(\mathbf{C}[\mathrm{G}]\) 的一组规范正交基。*

\BourbakiExercises{习题}

\begin{enumerate}
\def\labelenumi{\arabic{enumi})}
\item
  设 A 是一个交换环，G 是一个群，H 是 G 的一个有限指标的子群。假设 A 的元素 \((G:H)1_{A}\) 可逆。设 M 是一个 A{[}G{]}-模，N 是 M 的一个 A{[}G{]}-子模。证明：若 N 在 M 中容许一个补的 A{[}H{]}-子模，则它在 M 中容许一个补的 A{[}G{]}-子模（仿照 VIII, 第 401 页定理 1 的证明）。由此推出，一个作为 A{[}H{]}-模为投射模的 A{[}G{]}-模是投射模。
\item
  设 A 是一个交换环，G 是一个群。
\end{enumerate}

\begin{enumerate}
\def\labelenumi{\alph{enumi})}
\item
  假设群 G 的代数 A{[}G{]} 是一个绝对平坦环（VIII, 第 171 页，习题 27）。证明环 A 是绝对平坦的。设 H 是 G 的一个容许有限生成子集的子群；证明左理想 \(I_{H}\)（A{[}G{]} 的左理想，VIII, 第 22 页，习题 25）在 A{[}G{]} 中容许一个补，并由此推出 H 是有限的（注意，在相反的情形，\(I_{H}\) 的零化子将化为 (0)）。
\item
  保留 a) 的假设。设 g 是 G 的一个阶为有限数 n 的元素。证明 \(n1_{A}\) 在 A 中可逆（设 x 是 A 中满足 \((1-e_{g})x(1-e_{g})=(1-e_{g})\) 的元素；构造 A 的元素 y，使得 \(1-(1-e_{g})x=y(1+e_{g}+\cdots+e_{g^{n-1}})\)）。由此推出，G 的每个有限子群的阶都在 A 中可逆。
\item
  假设 A 是一个绝对平坦环，并且 G 的每个有限子集都生成一个基数在 A 中可逆的有限子群。证明环 A{[}G{]} 是绝对平坦的（归结为 G 有限的情形，并利用习题 1）。
\end{enumerate}

\begin{enumerate}
\def\labelenumi{\arabic{enumi})}
\setcounter{enumi}{2}
\tightlist
\item
  设 A 是一个交换环，G 是一个群。
\end{enumerate}

\begin{enumerate}
\def\labelenumi{\alph{enumi})}
\item
  假设群 G 的代数 A{[}G{]} 是半单的。证明 A 是半单的，G 是有限的，并且它的阶在 A 中可逆（利用习题 2 以及 VIII, 第 22 页的习题 25）。
\item
  假设代数 A{[}G{]} 同构于一个有限域族的乘积。证明 G 的每个子群都是正规的（利用 a)，并注意 A{[}G{]} 中的每个幂等元都是中心的）。
\end{enumerate}

在下面的习题 4 至 26 中，G 是一个有限群，K 是一个代数闭体；我们假设 G 的阶在 K 中可逆。

\begin{enumerate}
\def\labelenumi{\arabic{enumi})}
\setcounter{enumi}{3}
\item
  设 \((\mathrm{M},\pi)\) 是 G 的一个忠实的 \(^{(2)}\) 有限维表示，\((\mathrm{N},\rho)\) 是一个不可约表示；假设特征标 \(\chi_{\pi}\) 恰好取 s 个不同的值。证明存在整数 n \textless{} s，使得 K{[}G{]}-模 N 同构于 \(M^{\otimes n}\) 的一个子模（计算 \(\langle\chi_{\rho},\chi_{\pi}^{n}\rangle\)，并注意等式 \(\chi_{\pi}(g)=\dim M\) 等价于 g=1）。
\item
  设 \((\mathrm{M},\pi)\) 是 G 的一个表示。
\end{enumerate}

\begin{enumerate}
\def\labelenumi{\alph{enumi})}
\item
  设 H 是 G 的一个子群，S 是 G 中（模 H 的）左陪集的一个代表系，N 是 M 的一个在 H 下稳定的线性子空间。G 在 M 中的表示同构于由 H 在 N 中的表示诱导的表示，当且仅当 \(M = \bigoplus_{s \in S} sN\)。
\item
  假设 M 是子空间族 \((\mathrm{M}_{i})_{i\in\mathrm{I}}\) 的直和，并且存在 G 在 I 上的一个传递作用，使得 \(\pi(g)(\mathrm{M}_{i})\subset\mathrm{M}_{gi}\)（对 \(g\in G\)、\(i\in I\)）。设 \(o\in I\)，\(G_{o}\) 是它在 G 中的稳定子。证明 \(\pi\) 同构于由 \(G_{o}\) 在 \(M_{o}\) 中的表示诱导的表示。若 \(\pi\) 不可约，则 G 在 I 上的作用是传递的这一条件自动满足。
\end{enumerate}

\begin{enumerate}
\def\labelenumi{\arabic{enumi})}
\setcounter{enumi}{5}
\tightlist
\item
  设 X 是一个 G 作用于其上的有限集。我们考虑表示 \(\pi\)（G 在 \(K^{X}\) 中的表示），它满足 \(\pi(g)e_{x}=e_{gx}\)（对 \(g\in G\)、\(x\in X\)）（与 X 相伴的“置换表示”）。
\end{enumerate}

\begin{enumerate}
\def\labelenumi{\alph{enumi})}
\item
  对 \(g \in G\)，证明 \(\chi_{\pi}(g)\) 等于 X 中被 g 固定的点的个数。
\item
  证明单位表示在 \(\pi\) 中的重数等于 G 在 X 中的轨道的个数。
\item
  从现在起假设 G 的作用是传递的；取定 X 的一点 x，并用 H 表示它的稳定子。证明 \(\pi\) 同构于由 H 的单位表示诱导的表示，并且 \(\langle\chi_{\pi},\chi_{\pi}\rangle\) 等于 H 在 X 中的轨道的个数。
\item
  设 V 是 \(K^{X}\) 中由诸坐标之和为零的向量组成的子空间；它在 G 下稳定，我们用 \(\pi_{X}\) 表示由 \(\pi\) 得出的 G 在 V 中的表示。证明 \(\pi_{X}\) 不可约当且仅当 G 的作用是双传递的（I, §5, 第 143 页，习题 14）。
\end{enumerate}

\begin{enumerate}
\def\labelenumi{\arabic{enumi})}
\setcounter{enumi}{6}
\item
  对 G 的每个有限维表示 V，用 \(\varphi(\mathrm{V})\) 表示 V 中被 G 固定的子空间的维数；我们把 \(\varphi\) 线性延拓为 \(\mathrm{R}_{\mathrm{K}}(\mathrm{G})\) 上的一个 Z-线性型。设 \((\mathrm{V}_{\lambda})_{\lambda\in\widehat{\mathrm{G}}}\) 是 G 的不可约表示族，\(\omega\) 是元素 \(\sum_{\lambda\in\widehat{\mathrm{G}}}\left[\mathrm{V}_{\lambda}\right]\left[\mathrm{V}_{\lambda}^{*}\right]\)（\(\mathrm{R}_{\mathrm{K}}(\mathrm{G})\) 的元素）。证明公式 \(\varphi(\omega^{n})=\sum_{\mathrm{C}}d(\mathrm{C})^{n-1}\) 对每个整数 \(n\geqslant0\) 成立（应用 VIII, 第 415 页的命题 8 以及特征标的第二正交关系）。
\item
  设 H 与 F 是 G 的子群，S 是 G 的一个子集，使得 G 是诸双陪集 FsH 的不交并。设 \((N,\sigma)\) 是 H 的一个表示。对 \(s\in G\)，令 \(H_{s}=F\cap sHs^{-1}\)，并用 \(\sigma_{s}\) 表示 \(H_{s}\) 在 N 中的由 \(\sigma_{s}(x)=\sigma(s^{-1}xs)\) 定义的表示。
\end{enumerate}

\begin{enumerate}
\def\labelenumi{\alph{enumi})}
\item
  证明表示 \(\operatorname{Res}_{\mathrm{F}}^{\mathrm{G}}(\operatorname{Ind}_{\mathrm{H}}^{\mathrm{G}}(\sigma))\) 同构于诸表示 \(\operatorname{Ind}_{\mathrm{H}_{s}}^{\mathrm{F}}(\sigma_{s})\)（\(s \in S\)）的直和。（把表示 \(\operatorname{Ind}_{\mathrm{H}}^{\mathrm{G}}(\sigma)\) 记作 \((\mathrm{M}, \pi)\)。空间 M 可与诸 \(\pi(x)\mathrm{N}\)（\(x \in G/H\)）的直和等同。对 \(s \in S\)，设 \(\mathrm{M}(s)\) 是 M 中由诸 \(\pi(x)\mathrm{N}\)（\(x \in FsH\)）生成的子空间。注意 \(\mathrm{M}(s)\) 是诸子空间 \(\pi(xs)\mathrm{N}\)（\(x \in F/H_{s}\)）的直和，并由习题 5, b) 推出：作为 K{[}F{]}-模，\(\mathrm{M}(s)\) 同构于 \(\operatorname{Ind}_{\mathrm{H}_{s}}^{\mathrm{F}}(\sigma_{s})\)。
\item
  取 F = H（于是我们有 \(H_{s} = H \cap sHs^{-1}\)）。证明表示 \(\operatorname{Ind}_{\mathrm{H}}^{\mathrm{G}}(\sigma)\) 不可约，当且仅当 \(\sigma\) 不可约，并且对每个 \(s \in G - H\)，表示 \(\sigma_{s}\) 与 \(\operatorname{Res}_{\mathrm{H}_{s}}^{\mathrm{H}}(\sigma)\) 的支集（VIII, 第 66 页）不相交（“麦基判别法”：应用 a) 与弗罗贝尼乌斯互反律）。
\item
  假设 H 是正规的；这时 \(\operatorname{Ind}_{\mathrm{H}}^{\mathrm{G}}(\sigma)\) 不可约当且仅当 \(\sigma\) 不可约，且不与其共轭 \(\sigma \circ \operatorname{int}(g)\)（\(g \in G - H\)）中的任何一个同构。
\end{enumerate}

\begin{enumerate}
\def\labelenumi{\arabic{enumi})}
\setcounter{enumi}{8}
\tightlist
\item
  设 \((\mathrm{M},\pi)\) 是 G 的一个不可约表示，H 是 G 的一个正规子群。a) 设 \(M=\bigoplus_{i\in I}M_{i}\) 是 K{[}H{]}-模 M 分解为同型子模直和的分解。证明 G 在 M 上的作用传递地置换诸子模 \(M_{i}\)。由此推出 \((\mathrm{M},\pi)\) 同构于由 G 的一个包含 H 的子群诱导的表示（应用习题 5, b)）。
\end{enumerate}

\begin{enumerate}
\def\labelenumi{\alph{enumi})}
\setcounter{enumi}{1}
\tightlist
\item
  设 \(\mathcal{S}\) 是 K{[}H{]}-模 M 的支集；存在整数 e，使得 M（作为 K{[}H{]}-模）同构于 \(\sum_{\lambda\in \mathcal{S}}\lambda^{e}\)。
\end{enumerate}

\begin{enumerate}
\def\labelenumi{\arabic{enumi})}
\setcounter{enumi}{9}
\tightlist
\item
  设 \((\mathrm{M},\pi)\) 是 G 的一个不可约表示，H 是 G 的一个正规子群。a) 假设 H 等于 G 的中心。证明 \(\pi\) 的次数整除 H 在 G 中的指标。（对 \(m \geqslant 1\)，设 \(H_{m}\) 是 \(H^{m}\) 中由元素 \((h_{1},\ldots,h_{m})\)（满足 \(h_{1}\ldots h_{m}=1\) 的元素）组成的子群。注意表示 \(\pi^{\times m}:G^{m}\to\operatorname{End}_{K}(M^{\otimes m})\) 定义了 \(G^{m}/H_{m}\) 的一个不可约表示，并对它应用命题 10 的推论 2（VIII, 第 420 页）。）
\end{enumerate}

\begin{enumerate}
\def\labelenumi{\alph{enumi})}
\setcounter{enumi}{1}
\tightlist
\item
  假设 H 是交换的。证明 \(\pi\) 的次数整除 (G : H)。（利用习题 9, a)，并对 G 的基数作归纳，归结到 \(\pi\) 在 H 上的限制是同型的情形。证明 \(\pi(\mathrm{H})\) 在 \(\pi(\mathrm{G})\) 中是中心的，并应用 a)。）
\end{enumerate}

\begin{enumerate}
\def\labelenumi{\arabic{enumi})}
\setcounter{enumi}{10}
\tightlist
\item
  设 A 与 H 是 G 的子群。假设 A 是交换的且是正规的，并且 G 是 H 与 A 的半直积。群 H 作用在 \(\widehat{\mathrm{A}} = \mathrm{Hom}(\mathrm{A}, \mathrm{K}^{*})\) 上；取定代表系 \((\alpha_{i})_{i \in \widehat{\mathrm{A}}/\mathrm{H}}\)（H 在 \(\widehat{A}\) 中的轨道的代表系）。对 \(i \in \widehat{A}/H\)，设 \(H_{i}\) 是 \(\alpha_{i}\) 在 H 中的稳定子，并令 \(G_{i} = AH_{i}\)。对每个不可约表示 \(\sigma\)（\(H_{i}\) 的不可约表示），定义表示 \(\rho_{i,\sigma}\)（\(G_{i}\) 的一个表示）：令 \(\rho_{i,\sigma}(ah) = \alpha_{i}(a)\sigma(h)\)（对 \(a \in A\) 与 \(h \in H\)）。证明诸表示 \(\operatorname{Ind}_{G_{i}}^{G}(\rho_{i,\sigma})\)（\(i \in \widehat{A}/H\)，\(\sigma \in \widehat{H}_{i}\)）都是不可约的且两两不同构，并且用这种方法我们得到 G 的全部不可约表示（应用习题 9）。
\end{enumerate}

¶ 12) 设 H 是 G 的一个子群；假设对每个 \(g \in G - H\)，我们有等式 \(H \cap gHg^{-1} = \{1\}\)。

\begin{enumerate}
\def\labelenumi{\alph{enumi})}
\item
  设 u 是 H 上的一个中心函数。证明函数 \(\operatorname{Ind}_{\mathrm{H}}^{\mathrm{G}}(u)\) 在 \(H - \{1\}\) 上与 u 一致。若 v 是 H 上另一个满足 \(v(1) = 0\) 的中心函数，则我们有 \(\langle \operatorname{Ind}_{\mathrm{H}}^{\mathrm{G}}(u), \operatorname{Ind}_{\mathrm{H}}^{\mathrm{G}}(v) \rangle_{\mathrm{G}} = \langle u, v \rangle_{\mathrm{H}}\)。
\item
  设 \(\lambda \in \widehat{\mathrm{H}}\)；用 \(\chi_{\lambda}^{\prime}\) 表示 H 上的函数 \(\chi_{\lambda} - d_{\lambda}\)，用 \(\theta_{\lambda}\) 表示 G 上的函数 \(d_{\lambda} + \mathrm{Ind}_{\mathrm{H}}^{\mathrm{G}}(\chi_{\lambda}^{\prime})\)。证明 \(\theta_{\lambda}\) 是 G 的一个次数为 \(d_{\lambda}\)、在 H 上的限制同构于 \(\lambda\) 的不可约表示的特征标（利用 a) 计算 \(\langle\theta_{\lambda},\theta_{\lambda}\rangle\)、\(\langle\theta_{\lambda},1\rangle\) 与 \(\langle\mathrm{Res}_{\mathrm{H}}^{\mathrm{G}}(\theta_{\lambda}),\chi_{\lambda}\rangle\)）。
\item
  令 \(\theta = \sum_{\lambda} d_{\lambda} \theta_{\lambda}\)。证明：\(\theta(h) = 0\)（若 \(h \in \mathrm{H} - \{1\}\)），且 \(\theta(g) = \theta(1) = \mathrm{Card}(\mathrm{H})\)（若 \(g\) 不属于 \(\mathrm{H}\) 的任何共轭）。
\item
  设 \(L'\) 是 G 中不属于 H 的任何共轭的元素所成的集合，并令 \(L = L' \cup \{1\}\)。证明 L 是 G 的一个正规子群（由 c) 推出 L 是一个以 \(\theta\) 为特征标的表示的核）。
\item
  证明 G 是 H 与 L 的半直积（计算 L 的阶）。f) 设 \(h \in H - \{1\}\)。证明自同构 \(\operatorname{int}(h)\) 在 L 上的限制除 1 外没有不动点。由此推出 H 的阶整除 \(\operatorname{Card}(L) - 1\)。
\end{enumerate}

\(g\) ) 证明：若子群 H 的阶为偶数，则 L 是交换的（应用 I, §6, 第 145 页的习题 23, \(d\) )）。

\begin{enumerate}
\def\labelenumi{\arabic{enumi})}
\setcounter{enumi}{12}
\tightlist
\item
  设 \(\mathcal{H}\) 是 \(G\) 的子群所成的一个集合。
\end{enumerate}

\begin{enumerate}
\def\labelenumi{\alph{enumi})}
\tightlist
\item
  证明下列性质等价：
\end{enumerate}

\begin{enumerate}
\def\labelenumi{(\roman{enumi})}
\item
  属于 \(\mathcal{H}\) 的诸子群的共轭之并等于 G。
\item
  \(G\) 的每个特征标都是属于 \(\mathcal{H}\) 的子群的特征标所诱导的特征标的有理系数线性组合。
\end{enumerate}

（由 VIII, 第 412 页的命题 6 推出，性质 (ii) 等价于限制同态 \(\mathcal{Z}_{\mathrm{K}}(\mathrm{G})\to \bigoplus_{\mathrm{H}\in \mathcal{H}}\mathcal{Z}_{\mathrm{K}}(\mathrm{H}).\) 的单性。）

\begin{enumerate}
\def\labelenumi{\alph{enumi})}
\setcounter{enumi}{1}
\item
  设 \(\mathcal{C}\) 是 G 的循环子群所成的集合；它显然具有性质 (i)。对 G 的每个循环子群 C，用 \(\varphi_{\mathrm{C}}\) 表示 C 上的一个函数，它在 C 的生成元上取值 \(\operatorname{Card}(\mathrm{C})\)，在其他元素上取值 0。证明等式 \(\operatorname{Card}(\mathrm{G}) = \sum_{\mathrm{C} \in \mathcal{C}} \operatorname{Ind}_{\mathrm{C}}^{\mathrm{G}}(\varphi_{\mathrm{C}})\)。
\item
  设 \(C \in \mathcal{C}\)。证明函数 \(\varphi_{C}\) 属于子群 \(\mathrm{R}_{\mathrm{K}}(\mathrm{G})\)（\(\mathcal{Z}_{\mathrm{K}}(\mathrm{G})\) 的由诸特征标生成的子群）（对 C 的阶用归纳法，并利用 b)）。
\end{enumerate}

\(d\) ) 设 \(\chi\) 是 G 的一个不可约表示的特征标。证明等式 \(\chi = \mathrm{Card}(\mathrm{G})^{-1}\sum_{\mathrm{C}\in \mathcal{C}}\mathrm{Ind}_{\mathrm{C}}^{\mathrm{G}}(\varphi_{\mathrm{C}}\mathrm{Res}_{\mathrm{C}}^{\mathrm{G}}(\chi))\)，它给出了 \(a\) ) 中性质 (ii) 的一个显式公式。

\begin{enumerate}
\def\labelenumi{\arabic{enumi})}
\setcounter{enumi}{13}
\tightlist
\item
  若 G 的一个表示是它的子群的 1 次表示所诱导的表示的直和，则称它是单项表示。
\end{enumerate}

\begin{enumerate}
\def\labelenumi{\alph{enumi})}
\item
  假设 G 的真子群的表示都是单项的，并且 G 包含一个非中心的交换正规子群 H。证明 G 的每个忠实的不可约表示都是单项的（注意这样一个表示在 H 上的限制不是同型的，并应用习题 9, a)）。
\item
  由此推出，超可解群（I, §6, 第 146 页，习题 26），特别是幂零群的所有表示都是单项的。
\item
  取 G 为群 \(\mathbf{SL}(2,\mathbf{F}_{3})\)。它的中心 Z 的阶为 2，并且群 G/Z 同构于交错群 \(\mathfrak{A}_{4}\)（II, §10, 第 422 页，习题 14, g)）；特别地，G 是可解的。证明它的导群的指标为 3（参见 I, §5, 第 142 页，
\end{enumerate}

习题 10）。试由 VIII, 第 407 页的公式 (8) 推出 G 有一个次数为 2 的不可约表示，并证明该表示不是单项的。

¶ 15) 假设 G 的所有表示都是单项的；我们来证明 G 是可解的。对 G 的阶作归纳法，我们可以假设 G 的每个异于 G 的商群都是可解的。

\begin{enumerate}
\def\labelenumi{\alph{enumi})}
\tightlist
\item
  若 G 包含两个不同的非平凡极小正规子群 \(H_{1}\) 与 \(H_{2}\)，则 G 是可解的（考虑从 G 到 \(G/H_{1} \times G/H_{2}\) 的典范同态）。
\end{enumerate}

从现在起假设 G 有唯一的非平凡极小正规子群 H。

\begin{enumerate}
\def\labelenumi{\alph{enumi})}
\setcounter{enumi}{1}
\tightlist
\item
  证明：若 G 的一个表示在 H 上的限制是非平凡的，则它是忠实的。c) 设 \(\pi\) 是一个极小次数的忠实表示；存在 G 的子群 \(H_{i}\) 与表示 \(\sigma_{i}\)（\(H_{i}\) 的次数为 1 的表示），使得 \(\pi\) 同构于 \(\bigoplus\operatorname{Ind}_{\operatorname{H}_{i}}^{\operatorname{G}}(\sigma_{i})\)。令 \(\pi_{0}=\bigoplus\operatorname{Ind}_{\operatorname{H}_{i}}^{\operatorname{G}}(\varepsilon_{\operatorname{H}_{i}})\)，其中 \(\varepsilon_{H_{i}}\) 表示单位表示。证明 \(\pi_{0}\) 的核是交换的且非平凡（注意 \(\pi_{0}\) 是不可约的）；由此推出 G 是可解的。
\end{enumerate}

\begin{enumerate}
\def\labelenumi{\arabic{enumi})}
\setcounter{enumi}{15}
\tightlist
\item
  假设 K 的特征为零。我们说 K 的元素 x 在 Z 上是整的，如果它属于 K 的一个作为有限生成 Z-模的子环，或者等价地，如果 K 的子环 \(Z[x]\) 是一个有限生成的 Z-模 \(^{*}\)（参见 Comm. Alg., V, §1, No.~1, 第 303 页）。
\end{enumerate}

\begin{enumerate}
\def\labelenumi{\alph{enumi})}
\item
  证明 K 中在 Z 上是整的元素所成的集合是 K 的一个子环（对在 Z 上都是整的 x, y，考虑子环 \(Z[x, y]\)）。
\item
  元素 \(x\)（\(K\) 的元素）在 \(\mathbf{Z}\) 上是整的，当且仅当它是 \(\mathbf{Z}[X]\) 中某个首一多项式的根（当 \(x\) 是整的时候，仿照 VIII, 第 81 页引理 1 的证明，证明它满足形如 \(\det(xI_m - A) = 0\) 的方程，其中 \(A \in \mathbf{M}_m(\mathbf{Z})\)）。
\item
  若 x 在 Z 上是整的，则它的在 Q 上的共轭 \(x_{1},\ldots,x_{m}\) 也是整的；元素 \(x_{1}\cdots x_{m}\) 属于 Z（利用 VIII, 第 416 页的引理）。
\end{enumerate}

\begin{enumerate}
\def\labelenumi{\arabic{enumi})}
\setcounter{enumi}{16}
\tightlist
\item
  设 \(\mathrm{C}\) 是 \(\mathrm{G}\) 的一个共轭类，\(\pi\) 是 \(\mathrm{G}\) 的一个次数为 \(d\) 的不可约表示。
\end{enumerate}

\begin{enumerate}
\def\labelenumi{\alph{enumi})}
\item
  证明 K 的元素 \(d^{-1}\operatorname{Card}(\mathrm{C})\chi(\mathrm{C})\) 在 Z 上是整的（按 VIII, 第 415 页命题 9 的证明的方式推理）。
\item
  假设 d 与 \(\operatorname{Card}(\mathbf{C})\) 是互素的整数。证明元素 \(d^{-1}\chi(\mathbf{C})\) 在 Z 上是整的。
\end{enumerate}

* \(c\) ) 此外假设 \(\mathrm{K} = \mathbf{C}\)。证明 \(d^{-1}\chi (\mathrm{C})\) 在 \(\mathbf{Q}\) 上的所有共轭的绝对值都 \(\leqslant 1\)。由习题 16, \(c\) ) 推出 \(|\chi (\mathrm{C})|\) 或者为零或者等于 \(d\)。由此得出结论：若 \(\chi (\mathrm{C})\) 非零，则 \(\pi (x)\)（对每个 \(x\in \mathrm{C}\)）都是一个位似。特别地，若 \(\pi\) 是忠实的且 C 不化为单个元素，则我们有 \(\chi (\mathrm{C}) = 0.\ast\)

\begin{enumerate}
\def\labelenumi{\arabic{enumi})}
\setcounter{enumi}{17}
\tightlist
\item
  假设群 G 是单群但不是循环群。
\end{enumerate}

\begin{enumerate}
\def\labelenumi{\alph{enumi})}
\item
  证明不存在基数是素数幂且异于 \(\{1\}\) 的共轭类。（设 C 是一个基数为 \(p^{r}\)（\(r \geqslant 1\)）的共轭类。由习题 17, b) 推出，\(p^{-1}\chi_{\lambda}(1)\chi_{\lambda}(C)\)（当 \(\lambda\) 异于单位表示时）在 Z 上是整的，并用特征标的第二正交关系导出矛盾。）
\item
  证明 G 的阶至少被三个不同的素数整除（对 G 的一个西罗子群中的中心元素 \(x \neq 1\) 的共轭类应用 a)）。
\item
  证明阶至多被两个不同素数整除的有限群是可解的（“伯恩赛德定理”：对群的基数用归纳法，并利用 b)）。
\end{enumerate}

*19) 设 \(\pi\) 是 G 的一个次数 \(>1\) 的不可约复表示。证明存在 G 的一个元素 \(x\)，使得 \(\chi_{\pi}(x)=0\)。（设 \(\{\lambda_1,\ldots,\lambda_s\}\) 是 \(\pi\) 在 \(\widehat{G}\) 中的轨道（群 \(\operatorname{Gal}(\mathbf{C}/\mathbf{Q})\) 作用下的轨道），\(\chi_1,\ldots,\chi_s\) 是相应的特征标。设 \(g\in G\)；若 \(\chi(g)\neq 0\)，则仿照习题 17, \(c\) ) 的方式证明 \(|\chi_1(g)\ldots\chi_s(g)|\geqslant 1\)，从而 \(\sum_i|\chi_i(g)|^2\geqslant s\)（FRV, III, §1, No.~1, 第 93 页，命题 2）。若这对每个 \(g\in G\) 都成立，则由特征标的正交关系推出我们取等号，而取 \(g=1\) 就导出矛盾。）*

\begin{enumerate}
\def\labelenumi{\arabic{enumi})}
\setcounter{enumi}{19}
\tightlist
\item
  设 G 是一个有限群，H 是 G 的一个子群，\(\sigma\) 是 H 的一个有限维表示，\(\rho\) 是由 \(\sigma\) 诱导的 G 的表示。设 C 是 G 的一个共轭类；集合 \(C \cap H\) 是 H 的共轭类的一个有限族 \((D_{i})_{i \in I}\) 的不交并。
\end{enumerate}

\begin{enumerate}
\def\labelenumi{\alph{enumi})}
\item
  证明公式 \(\chi_{\rho}(\mathrm{C}) = \frac{(\mathrm{G:H})}{\mathrm{Card}(\mathrm{C})}\sum_{i\in \mathrm{I}}\mathrm{Card}(\mathrm{D}_i)\chi_\sigma (\mathrm{D}_i)\)。
\item
  若 \(\sigma\) 是平凡表示，则有 \(\chi_{\rho}(\mathrm{C})=(\mathrm{G}:\mathrm{H})\operatorname{Card}(\mathrm{C})^{-1}\operatorname{Card}(\mathrm{C}\cap\mathrm{H})\)。
\end{enumerate}

¶ 21) 设 \(n \in \mathbf{N}\)。用 \(\mathcal{D}_{n}\) 表示 \(\mathbf{N}^{n}\) 中由元素 \((p_{1},\ldots,p_{n})\)（满足 \(p_{1} \geqslant \cdots \geqslant p_{n}\) 的元素）组成的子集；赋予它字典序。对 \(\mathbf{p} \in \mathcal{D}_{n}\)，用 \(\mathbf{p}'\) 表示 \(\mathcal{D}_{n}\) 中由 \(p_{i}' = p_{i} + n - i\) 定义的元素。设 \(\mathbf{X} = (\mathrm{X}_{1},\ldots,\mathrm{X}_{n})\) 是一个变量序列；令 \(\Delta(\mathbf{X}) = \prod_{i<j}(\mathrm{X}_{i} - \mathrm{X}_{j})\)。

\begin{enumerate}
\def\labelenumi{\alph{enumi})}
\item
  设 \(\mathrm{P} \in \mathbf{Z}[\mathbf{X}]\) 是一个反对称多项式，即满足关系 \(\mathrm{P}(\mathrm{X}_{\sigma(1)}, \ldots, \mathrm{X}_{\sigma(n)}) = \varepsilon(\sigma)\mathrm{P}(\mathrm{X}_1, \ldots, \mathrm{X}_n)\)（对每个元素 \(\sigma \in \mathfrak{S}_n\)）的多项式。证明存在对称多项式 \(\mathrm{Q} \in \mathbf{Z}[\mathbf{X}]\)，使得 \(\mathrm{P} = \Delta\mathrm{Q}\)。
\item
  对每个 \(\mathbf{p} \in \mathcal{D}_n\)，令 \(S_{\mathbf{p}}(\mathbf{X}) = \Delta(\mathbf{X})^{-1} \det(X_j^{p_i'})_{1 \leqslant i,j \leqslant n}\)。证明 \(S_{\mathbf{p}}\) 是一个次数为 \(\sum p_i\)、具有整系数的齐次对称多项式。
\item
  若 \(\mathbf{q} = (q_i)_{1\leqslant i\leqslant n}\) 是 \(\mathcal{D}_n\) 的一个元素，而 \(P\) 是 \(\mathbf{Q}[\mathbf{X}]\) 的一个元素，则用 \([P]_{\mathbf{q}}\) 表示单项式 \(\mathbf{x}^{\mathbf{q}}\) 在 \(P\) 中的系数。令 \(K_{\mathbf{pq}} = [S_{\mathbf{p}}]_{\mathbf{q}}\)。证明我们有 \(K_{\mathbf{pp}} = 1\) 与 \(K_{\mathbf{pq}} = 0\)（对 \(\mathbf{q} > \mathbf{p}\)）（按字典序展开等式 \(\det(X_j^{p_i'}) = S_{\mathbf{p}}(\mathbf{X})\Delta(\mathbf{X})\)）。由此推出，诸多项式 \(S_{p}\)（\(p \in \mathcal{D}_{n}\)）构成 Z{[}X{]} 中由对称多项式组成的 Z-子模的一组基。
\item
  设 \(\mathbf{Y} = (\mathrm{Y}_1, \ldots, \mathrm{Y}_n)\) 是一个变量序列。证明形式幂级数的等式
\end{enumerate}

\[
\prod_ {1 \leqslant i, j \leqslant n} (1 - \mathrm{X} _ {i} \mathrm{Y} _ {j}) ^ {- 1} = \sum_ {\mathbf {p} \in \mathcal {D} _ {n}} \mathrm{S} _ {\mathbf {p}} (\mathbf {X})   \mathrm{S} _ {\mathbf {p}} (\mathbf {Y})
\]

（利用等式 \(\det\big((1 - X_iY_j)^{-1}\big)_{1\leqslant i,j\leqslant n} = \Delta(\mathbf{X})\Delta(\mathbf{Y})\prod_{1\leqslant i,j\leqslant n}(1 - X_iY_j)^{-1}\)（参见 III, §8, 第 639 页，习题 5），把 \((1 - T)^{-1}\) 作为形式幂级数展开以计算左边）。

\begin{enumerate}
\def\labelenumi{\alph{enumi})}
\setcounter{enumi}{4}
\tightlist
\item
  设 \(\mathrm{P} \in \mathbf{Q}[\mathbf{X}]\) 是一个对称多项式。对每个 \(\mathbf{p} \in \mathcal{D}_n\)，令 \(\omega_{\mathbf{p}}(\mathrm{P}) = [\mathrm{P}\Delta]_{\mathbf{p}'}\)。证明等式 \(\omega_{\mathbf{p}}(\mathrm{S}_{\mathbf{q}}) = \delta_{\mathbf{pq}}\)。由此推出，对每个 \(\mathbf{q} \in \mathcal{D}_n\)，我们有 \([\mathrm{P}]_{\mathbf{q}} = \sum_{\mathbf{p} \in \mathcal{D}_n} \mathrm{K}_{\mathbf{pq}} \omega_{\mathbf{p}}(\mathrm{P})\)。
\end{enumerate}

\(f)\) 对 \(\boldsymbol{\alpha}\in\mathbf{N}^{n}\)，令 \(f^{\boldsymbol{\alpha}}(\mathbf{X})=\prod_{j=1}^{n}(\sum_{i}\mathrm{X}_{i}^{j})^{\alpha_{j}}\) 与 \(c_{\boldsymbol{\alpha}}=\prod_{j=1}^{n}\frac{1}{j^{\alpha_{j}}\alpha_{j}!}\)。证明我们有 \(\sum_{\boldsymbol{\alpha}\in\mathbf{N}^{n}}c_{\boldsymbol{\alpha}}\omega_{\mathbf{p}}(f^{\boldsymbol{\alpha}})\omega_{\mathbf{q}}(f^{\boldsymbol{\alpha}})=\delta_{\mathbf{pq}}\)（对 \(\mathbf{p},\mathbf{q}\)（\(\mathcal{D}_{n}\) 中的元素））（注意我们有

\[
\prod_ {1 \leqslant i, j \leqslant n} (1 - \mathrm{X} _ {i} \mathrm{Y} _ {j}) ^ {- 1} = \prod_ {k = 0} ^ {\infty} \exp \bigl (\frac {1}{k} (\sum_ {i} \mathrm{X} _ {i} ^ {k}) (\sum_ {j} \mathrm{Y} _ {j} ^ {k}) \bigr) = \sum_ {\boldsymbol {\alpha} \in \mathbf {N} ^ {n}} c _ {\boldsymbol {\alpha}} f ^ {\boldsymbol {\alpha}} (\mathbf {X}) f ^ {\boldsymbol {\alpha}} (\mathbf {Y}) \bigr).
\]

¶ 22) 设 n 是一个自然数。用 \(\mathcal{S}_{n}\) 表示和为 n 的严格正整数的递减有限序列所成的集合。设 \(\mathbf{p} = (p_{i})_{1 \leqslant i \leqslant d}\) 是 \(\mathcal{S}_{n}\) 的一个元素。对 \(1 \leqslant i \leqslant d\)，用 \(P_{i}\) 表示满足 \(p_{1} + \cdots + p_{i-1} < r \leqslant p_{1} + \cdots + p_{i}\) 的整数 r 所成的集合；对 \(1 \leqslant j \leqslant p_{1}\)，用 \(P_{j}'\) 表示形如 \(p_{1} + \cdots + p_{i-1} + j\)（其中 \(1 \leqslant i \leqslant d\) 且 \(p_{i} \geqslant j\)）的整数所成的集合。子集 \(P_{i}\) 一方面、子集 \(P_{j}'\) 另一方面，构成 N 中区间 {[}1, n{]} 的一个分拆。

（我们常用一个图来表示这一情形，该图称为杨图：在下面的例子中，我们有 n = 9 与 \(\mathbf{p} = (4, 2, 2, 1)\)。集合 \(P_{i}\) 由第 i 行的元素组成，而 \(P_{j}^{\prime}\) 由第 j 列的元素组成。）

1

2

3

4

5

6

7

8

9

设 \(\mathcal{P}\)（相应地，\(\mathcal{P}'\)）是 \(\mathfrak{S}_n\) 中由这样的置换 \(\sigma\) 组成的子群：\(\sigma(\mathrm{P}_i) \subset \mathrm{P}_i\)（对每个 \(i\)）且（相应地）\(\sigma(\mathrm{P}_j') \subset \mathrm{P}_j'\)（对每个 \(j\)）。

\begin{enumerate}
\def\labelenumi{\alph{enumi})}
\item
  证明我们有 \(\mathcal{P} \cap \mathcal{P}' = \{1\}\)。
\item
  设 \(\sigma \in \mathfrak{S}_n\)。证明：若 \(\mathrm{Card}(\mathrm{P}_i \cap \sigma(\mathrm{P}_j')) \leqslant 1\)（对每个 \(i\) 与每个 \(j\)），则 \(\sigma\) 属于 \(\mathcal{PP}'\)。由此推出，若 \(\sigma \notin \mathcal{PP}'\)，则子群 \(\mathcal{P} \cap \sigma \mathcal{P}'\sigma^{-1}\) 包含一个对换。
\item
  在群代数 \(\mathbf{Z}[\mathfrak{S}_n]\) 中，令
\end{enumerate}

\[
a _ {\mathbf {p}} = \sum_ {\sigma \in \mathcal {P}} e _ {\sigma}, b _ {\mathbf {p}} = \sum_ {\tau \in \mathcal {P} ^ {\prime}} \varepsilon (\tau) e _ {\tau}, c _ {\mathbf {p}} = a _ {\mathbf {p}} b _ {\mathbf {p}}.
\]

证明 \(c_{p}\) 满足 \(e_{\sigma}c_{p}e_{\sigma'}=\varepsilon(\sigma')c_{p}\)（对每个 \(\sigma\in \mathcal{P}\) 与每个 \(\sigma'\in \mathcal{P}'\)），并且具有这一性质的每个元素都属于 \(Zc_{p}\)（把这样一个元素写成 \(\sum n_{\sigma}e_{\sigma}\)，并由 b) 推出 \(n_{\sigma}\) 为零（当 \(\sigma\notin \mathcal{P}\mathcal{P}'\)））。由此推出，对 \(Z[\mathfrak{S}_n]\) 的每个元素 x，元素 \(c_{p}xc_{p}\) 属于 \(Zc_{p}\)。

\(d\) ) 用 A 表示 \(\mathbf{Q}\)-代数 \(\mathbf{Q}[\mathfrak{S}_n]\)。证明 A 的理想 \(V_{\mathbf{p}} = Ac_{\mathbf{p}}\) 是一个绝对单的 A-模。（注意我们有 \(c_{\mathbf{p}}V_{\mathbf{p}} \subset \mathbf{Q}c_{\mathbf{p}}\)，由 \(c\)。）证明：若 \(\mathfrak{a}\) 是一个真包含于 \(V_{\mathbf{p}}\) 之中的左理想，则我们有 \(c_{\mathbf{p}}\mathfrak{a} = 0\)，从而 \(\mathfrak{a}^2 = 0\)，最终 \(\mathfrak{a} = 0\)。）

\begin{enumerate}
\def\labelenumi{\alph{enumi})}
\setcounter{enumi}{4}
\tightlist
\item
  证明我们有 \(c_{\mathbf{p}}^{2} = (n!/[V_{\mathbf{p}} : \mathbf{Q}]) c_{\mathbf{p}}\)（计算 A 的自同态 \(x \mapsto xc_{p}\) 的迹）。
\end{enumerate}

¶ 23) 沿用上一习题的记号。设 \(\mathbf{p} = (p_i)\) 与 \(\mathbf{q} = (q_j)\) 是 \(S_n\) 的两个元素。

\begin{enumerate}
\def\labelenumi{\alph{enumi})}
\item
  假设按字典序有 \(\mathbf{p} > \mathbf{q}\)。证明我们有 \(a_{\mathbf{p}}xb_{\mathbf{q}} = 0\)（对每个 \(x \in \mathrm{A}\)）。（归结为 \(x = 1\) 的情形。注意若 \((\mathrm{P}_i), (\mathrm{P}_j')\) 表示与 \(\mathbf{p}\) 相伴的分拆，而 \((\mathrm{Q}_k), (\mathrm{Q}_l')\) 表示与 \(\mathbf{q}\) 相伴的分拆，则记 \(s\) 为最小的整数 \(n\)，使得 \(p_n \neq q_n\)，我们有 \(\operatorname{Card}(\mathrm{P}_s \cap \mathrm{Q}_1') \geqslant 2\)；构造一个对换 \(\tau\)，使得 \(a_{\mathbf{p}}\tau = a_{\mathbf{p}}\) 且 \(\tau b_{\mathbf{q}} = -b_{\mathbf{q}}\)。）
\item
  由此推出，若 \(p \neq q\)，则 A-模 \(V_{p}\) 与 \(V_{q}\) 不同构（应用 a)，并利用把 \(e_{\sigma}\) 映为 \(e_{\sigma^{-1}}\) 的 A 的反自同构）。
\item
  证明把序列 \(\mathbf{p}\) 映为表示 \((\mathrm{V}_{\mathbf{p}})_{(\mathrm{K})}\) 的映射定义了一个从 \(\mathcal{S}_n\) 到集合 \(\mathcal{S}_{\mathrm{K}}(\mathfrak{S}_n)\)（单 \(\mathrm{K}[\mathfrak{S}_n]\)-模的类所成的集合）的双射。d) 对 \(\mathbf{p} \in \mathcal{S}_n\)，用 \(\overline{\mathbf{p}}\) 表示由 \(\overline{p}_j = \mathrm{Card}(\mathrm{P}_j')\) 定义的序列（习题 22）。证明映射 \(\mathbf{p} \mapsto \overline{\mathbf{p}}\) 是 \(\mathcal{S}_n\) 的一个对合；A-模 \(\mathrm{V}_{\overline{\mathbf{p}}}\) 同构于 \(\mathrm{V}_{\mathbf{p}} \otimes_{\mathbf{Q}} \mathbf{Q}_{\varepsilon}\)，其中 \(\mathbf{Q}_{\varepsilon}\) 表示 \(\mathbf{Q}\) 上与符号相伴的一维表示。
\item
  描述与序列 \((1, \ldots, 1)\) 和 \((n)\) 相对应的表示，以及与序列 \((n - 1, 1)\) 相对应的表示。
\end{enumerate}

¶ 24) 沿用习题 21 至 23 的记号。设 \(\mathbf{p} = (p_i)_{1 \leqslant i \leqslant d} \in \mathcal{S}_n\)；用 \(\mathbf{p}\) 表示元素 \((p_1, \ldots, p_d, 0, \ldots, 0)\)（\(\mathcal{D}_n\) 的元素）。若 \(\boldsymbol{\alpha} = (\alpha_1, \ldots, \alpha_n)\) 是 \(\mathbf{N}^n\) 的满足 \(\sum_i i\alpha_i = n\) 的一个元素，则用 \(C_{\boldsymbol{\alpha}}\) 表示 \(\mathfrak{S}_n\) 的由如下元素组成的共轭类：这些元素是一族支集两两不交的轮换的乘积，其中 \(\alpha_2\) 个轮换的阶为 2，\(\alpha_3\) 个轮换的阶为 3，等等。

\begin{enumerate}
\def\labelenumi{\alph{enumi})}
\item
  证明表示 \(\rho_{\mathbf{p}}\)（\(\mathfrak{S}_n\) 在 A-模 \(\mathrm{Aa}_{\mathbf{p}}\) 中的表示）是由 \(\mathcal{P}\) 的平凡表示诱导的表示。利用习题 20 推出公式 \(\chi_{\rho_{\mathbf{p}}}(\mathrm{C}_{\alpha}) = [f^{\alpha}]_{\mathbf{p}}\)。
\item
  对任意 \(\mathbf{q} \in \mathcal{S}_n\)，用 \(\lambda_{\mathbf{q}}\) 表示 \(\mathfrak{S}_n\) 在 A-模 \(V_{\mathbf{q}}\) 中的表示。证明存在正整数 \(n_{\mathbf{pq}}\)，使得我们有 \(\chi_{\rho_{\mathbf{p}}} = \sum_{\mathbf{q} \in \mathcal{S}_n} n_{\mathbf{pq}} \chi_{\lambda_{\mathbf{q}}}\)（其中 \(n_{\mathbf{pp}} > 0\)）。
\item
  设 \(\eta_{\mathbf{p}}\) 是 \(\mathfrak{S}_n\) 上的取值 \(\omega_{\mathbf{p}}(f^{\alpha})\)（在类 \(\mathrm{C}_{\alpha}\) 的元素上）的中心函数。证明 \(\eta_{\mathbf{p}}\) 是诸特征标 \(\chi_{\lambda_{\mathbf{q}}}\)（\(\mathbf{q} \in \mathcal{S}_n\)）的具有整系数的线性组合（利用习题 21, c) 与 e)）。由此推出，\(\eta_{\mathbf{p}}\) 相差一个符号等于某个特征标 \(\chi_{\lambda_{\mathbf{q}}}\)。（利用习题 21, f) 计算 \(\langle \eta_{\mathbf{p}}, \eta_{\mathbf{p}} \rangle\)。）d) 推出等式 \(\chi_{\lambda_{\mathbf{p}}}(\mathrm{C}_{\alpha}) = [\Delta f^{\alpha}]_{\mathbf{p}'}\)（“弗罗贝尼乌斯公式”）。
\item
  证明不可约表示 \(\lambda_{p}\) 在表示 \(\rho_{q}\) 中的重数等于 \(K_{pq}\)。
\end{enumerate}

\(f\) ) 证明 \(\lambda_{\mathbf{p}}\) 的次数为 \(\frac{n!}{\mathbf{p}^{\prime}!}\Delta (\mathbf{p}^{\prime})\)（展开多项式 \(\Delta (\mathbf{X})(\sum \mathrm{X}_i)^n\)）。

¶ 25) 考虑表示 \(\pi_{\mathbf{n}}\)（\(\mathfrak{S}_n\) 的与 \(\mathfrak{S}_n\) 在集合 \(\mathbf{n} = [1, n]\) 上的作用相伴的表示）（习题 6）。设 \(k\) 是一个正整数，\(\chi_k\) 是表示 \(\wedge^k \pi_{\mathbf{n}}\) 的特征标。

\begin{enumerate}
\def\labelenumi{\alph{enumi})}
\item
  设 \(\sigma \in \mathfrak{S}_n\)。证明公式 \(\chi_k(\sigma) = \sum \varepsilon (\sigma |S)\)，其中求和取遍诸子集 \(S\)（\(\mathbf{n}\) 的具有 \(k\) 个元素且满足 \(\sigma (S) = S\) 的子集）所成的集合。
\item
  证明 \(\wedge^{k}\pi_{n}\) 是 \(\mathfrak{S}_n\) 的与元素 \((n-k,1,\ldots,1)\)（\(\mathcal{S}_n\) 的元素）相伴的一个不可约表示（利用弗罗贝尼乌斯公式计算这个表示的特征标，并注意和式
\end{enumerate}

\[
\Delta (\mathbf {X}) f ^ {\boldsymbol {\alpha}} (\mathbf {X}) = \sum_ {\sigma \in \mathfrak {S} _ {n}} \varepsilon (\sigma) X _ {1} ^ {\sigma (n) - 1} \dots X _ {n} ^ {\sigma (1) - 1} f ^ {\boldsymbol {\alpha}} (\mathbf {X})
\]

中 \(X^{p'}\) 的系数非零的那些项对应于这样的置换 \(\sigma\)：对每个 i 满足 \(\sigma(i) \leqslant i + 1\)，且满足 \(\sigma(i) = i\)（对 \(1 \leqslant i < n - k\)）。

¶ 26) 设 q 是一个素数的幂，F 是有 q 个元素的有限域；取 G 为群 \(\mathbf{GL}(2, \mathbf{F})\)。

\begin{enumerate}
\def\labelenumi{\alph{enumi})}
\item
  描述 G 的共轭类（区分四种类型的类，它们分别含有 1、\(q^{2}-1\)、\(q^{2}+q\) 与 \(q^{2}-q\) 个元素）。
\item
  群 G 作用在 F 上的射影直线 P 上；考虑 G 的表示 \(\pi_{P}\)（习题 6）。对任意群同态 \(\lambda : F^{*} \to K^{*}\)，用 \(\delta_{\lambda}\) 表示与同态 \(\lambda \circ det\)（从 G 到 \(K^{*}\) 的同态）相伴的次数为 1 的表示，并令 \(\pi_{\lambda} = \pi_{P} \otimes \delta_{\lambda}\)。证明 \(\pi_{\lambda}\) 不可约，并计算它的特征标。
\item
  设 B 是由上三角矩阵组成的 G 的子群。对 \(\lambda,\mu\)（\(\mathrm{Hom}(\mathrm{F}^{*},\mathrm{K}^{*})\) 中的元素），用 \(\beta_{\lambda,\mu}\) 表示由 \(\beta_{\lambda,\mu}((a_{ij}))=\lambda(a_{11})\mu(a_{22})\) 定义的 B 的次数为 1 的表示，并令 \(\pi_{\lambda,\mu}=\mathrm{Ind}_{\mathrm{B}}^{\mathrm{G}}(\beta_{\lambda,\mu})\)。证明 \(\pi_{\lambda,\mu}\) 当 \(\lambda\neq\mu\) 时不可约，而 \(\pi_{\lambda,\lambda}\) 同构于 \(\pi_{P}\oplus\delta_{\lambda}\)（应用习题 8）。表示 \(\pi_{\lambda,\mu}\) 与 \(\pi_{\lambda',\mu'}\) 同构，当且仅当子集 \(\{\lambda,\mu\}\) 与 \(\{\lambda',\mu'\}\)（\(\mathrm{Hom}(\mathrm{F}^{*},\mathrm{K}^{*})\) 的子集）相等。
\end{enumerate}

\(d\) ) 设 \(\mathrm{F}^{\prime}\) 是 F 的一个次数为 2 的扩张；取定 \(\mathrm{F}'\) 在 F 上的一组基，就可把 \(\mathrm{F}^{\prime *}\) 与 \(\mathrm{G} = \mathrm{Aut}_{\mathrm{F}}(\mathrm{F}')\) 的一个子群等同起来。设 \(\varphi \in \mathrm{Hom}(\mathrm{F}^{\prime *},\mathrm{K}^{*})\)；计算 \(\mathrm{Ind}_{\mathrm{F}^{\prime *}}^{\mathrm{G}}(\varphi)\)，并证明它等于 \(\mathrm{Ind}_{\mathrm{F}^{\prime *}}^{\mathrm{G}}(\varphi^q)\)。用 \(\lambda\) 表示 \(\varphi\) 在 \(\mathrm{F}^*\) 上的限制，并令 \(\chi_{\varphi} = \chi_{\pi_{\mathrm{P}}\otimes \pi_{\lambda ,1}} - \chi_{\pi_{\lambda ,1}} - \mathrm{Ind}_{\mathrm{F}^{\prime *}}^{\mathrm{G}}(\varphi)\)。证明：若 \(\varphi \neq \varphi^q\)，则我们有 \(\langle \chi_{\varphi},\chi_{\varphi}\rangle = 1\) 与 \(\chi_{\varphi}(1) = q - 1\)，从而 \(\chi_{\varphi}\) 是一个不可约表示 \(\pi_{\varphi}\)（维数为 \(q - 1\) 的表示）的特征标。证明用这种方法我们得到 \(\frac{1}{2} q(q - 1)\) 个两两不同构的不可约表示。

\begin{enumerate}
\def\labelenumi{\alph{enumi})}
\setcounter{enumi}{4}
\tightlist
\item
  证明诸不可约表示 \(\pi_{\lambda}\)（\(\lambda\in\mathrm{Hom}(\mathrm{F}^{*},\mathrm{K}^{*})\)）、\(\pi_{\lambda,\mu}\)（\(\{\lambda,\mu\}\) 是 \(\mathrm{Hom}(\mathrm{F}^{*},\mathrm{K}^{*})\) 的一个二元素子集）与 \(\pi_{\varphi}\)（\(\varphi\in\mathrm{Hom}(\mathrm{F}^{\prime*},\mathrm{K}^{*})\) 且 \(\varphi\neq\varphi^{q}\)）两两不同构，并且 G 的每个不可约表示都同构于它们之一。
\end{enumerate}

\begin{enumerate}
\def\labelenumi{\arabic{enumi})}
\setcounter{enumi}{26}
\tightlist
\item
  设 F 是一个特征为 p \textgreater{} 0 的代数闭体，G 是一个有限群。
\end{enumerate}

\begin{enumerate}
\def\labelenumi{\alph{enumi})}
\item
  设 \(x \in G\)。证明存在唯一的分解 \(x = x_{u}x_{s}\)，其中 \(x_{u}\) 的阶是 p 的幂，\(x_{s}\) 的阶与 p 互素，且 \(x_{u}\) 与 \(x_{s}\) 交换。元素 \(x_{u}\) 与 \(x_{s}\) 都是 x 的幂。
\item
  设 \(\pi\) 是 \(G\) 在 \(F\) 上的一个有限次数线性表示。证明我们有 \(\chi_{\pi}(x) = \chi_{\pi}(x_s)\)（对每个 \(x \in G\)）。
\item
  若 G 的元素 x 的阶与 p 互素，则称它是 p-正则的；用 \(G^{reg}\) 表示 G 的 p-正则元素所成的集合，用 \(\mathcal{Z}_{\mathrm{F}}(\mathrm{G}^{\mathrm{reg}})\) 表示从 \(G^{reg}\) 到 F 的中心函数所成的空间。由同态 \(\mathrm{R}_{\mathrm{F}}(\mathrm{G}) \to \mathcal{Z}_{\mathrm{F}}(\mathrm{G}^{\mathrm{reg}})\)（把表示 \(\pi\) 的类映为 \(\chi_{\pi}\) 在 \(G^{reg}\) 上的限制的同态），我们导出一个环同态 \(\psi : F \otimes_{\mathbf{Z}} R_{\mathrm{F}}(G) \to \mathcal{Z}_{\mathrm{F}}(\mathrm{G}^{\mathrm{reg}})\)。证明 \(\psi\) 是单射（利用 b) 以及 VIII, 第 384 页的推论）。
\item
  设 x 是 G 的一个 p-正则元素，\(Z(x)\) 是它的交换子，P 是 \(Z(x)\) 的一个西罗 p-子群，\(H_{x}\) 是由 P 与 x 生成的 G 的子群。证明 \(H_{x}\) 是 P 与由 x 生成的 G 的子群的直积，并且它的共轭类由 x 的共轭类完全确定。
\item
  构造次数为 1 的表示 \(\sigma_{1},\ldots,\sigma_{m}\)（\(H_{x}\) 的表示）与 F 的元素 \(a_{1},\ldots,a_{m}\)，使得函数 \(\sum a_{i}\chi_{\sigma_{i}}\) 在 x 处取值 1，在 x 的其他幂处取值 0（先构造由 x 生成的子群的表示）。令 \(\rho_{i}=\operatorname{Ind}_{H_{x}}^{G}\pi_{i}\)（对 \(i=1,\ldots,m\)），并令 \(f=\sum_{i}a_{i}\chi_{\rho_{i}}\)。证明 \(f(x)\) 等于 \((\mathrm{Z}(x):\mathrm{H}_{x})1_{\mathrm{F}}\)——它在 F 中非零，并且 \(f(y)\) 对 G 的每个不与 x 共轭的 p-正则元素 y 都为零。
\end{enumerate}

\(f\) ) 由 \(e\) 推出，同态 \(\psi : \mathrm{F} \otimes_{\mathbf{Z}} \mathrm{R}_{\mathrm{F}}(\mathrm{G}) \to \mathcal{Z}_{\mathrm{F}}(\mathrm{G}^{\mathrm{reg}})\) 是双射。特别地，单 F{[}G{]}-模的同构类的个数等于 G 的 \(p\)-正则元素的共轭类的个数。

\begin{enumerate}
\def\labelenumi{\arabic{enumi})}
\setcounter{enumi}{27}
\tightlist
\item
  沿用习题 27 的假设。设 \(m\) 是 \(p\)-正则元素（\(G\) 的元素）的阶的最小公倍数。设 \(F_0\) 是 \(F\) 的由 \(m\) 次单位根生成的子域；它是一个有限域。
\end{enumerate}

\begin{enumerate}
\def\labelenumi{\alph{enumi})}
\item
  证明 G 在 F 上的任何有限次数表示的特征标的所有值都属于 \(F_{0}\)。
\item
  证明 \(F_{0}[G]\) 关于其根的商环是 \(F_{0}\) 上矩阵代数的一个乘积（利用 a) 与韦德伯恩定理）。由此推出，典范同态 \(\mathrm{R}_{\mathrm{F}_{0}}(\mathrm{G}) \to \mathrm{R}_{\mathrm{F}}(\mathrm{G})\) 是双射。
\end{enumerate}

\begin{enumerate}
\def\labelenumi{\arabic{enumi})}
\setcounter{enumi}{28}
\tightlist
\item
  沿用习题 27 与 28 的假设和记号。取定一个本原 \(m\) 次单位根 \(\zeta\)。设 \(\mathcal{O}_m\) 是环 \(\mathbf{Z}[\mathrm{X}] / (\Phi_m)\)（VIII, 第 415 页）；X 在 \(\mathcal{O}_m\) 中的类记作 \(x\)。设 \(\varphi : \mathcal{O}_m \to \mathrm{F}\) 是把 \(x\) 映为 \(\zeta\) 的同态。它诱导出从 \(\mathcal{O}_m\) 中由 \(x\) 生成的子群到群 \(\mu_m(\mathrm{F})\) 的一个同构；用 \(\tau\) 表示逆同构。
\end{enumerate}

\begin{enumerate}
\def\labelenumi{\alph{enumi})}
\item
  设 \(\pi\) 是 G 在 F 上的一个线性表示，维数 \(d\) 有限。对 \(g \in \mathrm{G}\)，设 \(\mathrm{P}_g(\mathrm{X}) = \prod_{i=1}^d (\mathrm{X} - \lambda_i)\) 是 \(\pi(g)\) 的特征多项式；用 \(\beta_{\pi}(g)\) 表示元素 \(\sum_{i} \tau(\lambda_i)\)（\(\mathcal{O}_m\) 的元素）。这样得到的函数 \(\beta_{\pi}: \mathrm{G} \to \mathcal{O}_m\) 称为 \(\pi\) 的布饶尔特征标；我们有 \(\varphi \circ \beta_{\pi} = \chi_{\pi}\) 与 \(\beta_{\pi}(x) = \beta_{\pi}(x_s)\)（对每个 \(x\)），从而 \(\beta_{\pi}\) 由它在 \(\mathrm{G}^{\mathrm{reg}}\) 上的限制确定。
\item
  设 \(\widehat{\mathbf{G}}\) 是单 F{[}G{]}-模的同构类所成的集合。证明诸函数 \((\beta_{\lambda})_{\lambda \in \widehat{\mathbf{G}}}\) 在 \(\mathbf{Z}\) 上线性无关。（若存在一个非平凡关系 \(\sum n_{\lambda} \beta_{\lambda} = 0\)，则可以假设某个 \(n_{\lambda}\) 与 \(p\) 互素，对它应用 \(\varphi\) 并利用习题 27, f) 得到矛盾。）
\item
  设 \(\mathcal{B}\) 是从 \(\mathrm{G}^{\mathrm{reg}}\) 到 \(\mathcal{O}_m\) 的函数所成的集合；同态 \(\mathrm{R}_{\mathrm{F}}(\mathrm{G}) \to \mathcal{B}\)（把 \(\pi\) 的类映为 \(\beta_{\pi}\) 的同态）是单射。由此推出，两个半单 F{[}G{]}-模同构，当且仅当它们的布饶尔特征标相等。
\end{enumerate}

\begin{enumerate}
\def\labelenumi{\arabic{enumi})}
\setcounter{enumi}{29}
\tightlist
\item
  设 \(\mathrm{F}\) 是一个特征为 \(p > 0\) 的代数闭体，\(\mathbf{F}_p\) 是 \(\mathrm{F}\) 的素子域，并设 \(\mathrm{G} = \mathbf{SL}_2(\mathbf{F}_p)\)。用 \(\mathrm{V}\) 表示 F-向量空间 \(\mathrm{F} \otimes_{\mathbf{F}_p} \mathbf{F}_p^2\)，它带有 \(\mathrm{G}\) 的一个作用，该作用由 \(\mathbf{F}_p^2\) 上的作用导出。
\end{enumerate}

\begin{enumerate}
\def\labelenumi{\alph{enumi})}
\item
  证明 F{[}G{]}-模 \(\mathbf{S}^{d}(V)\) 是单模（对 \(0 \leqslant d < p\)）。（设 W 是 F{[}G{]}-模 \(\mathbf{S}^{d}(V)\) 的一个非零子模，\((e, f)\) 是 V 的一组基。证明 \(e^{d}\) 属于 W（利用矩阵 \(\begin{pmatrix}1 & 1 \\ 0 & 1\end{pmatrix}\)），然后证明 W 等于 \(\mathbf{S}^{d}(V)\)（利用矩阵 \(\begin{pmatrix}1 & 0 \\ 1 & 1\end{pmatrix}\)）。）b) 证明 F{[}G{]}-模 \(\mathbf{S}^{p}(V)\) 不是单模（考虑 V 在映射 \(v \mapsto v^{p}\) 下的像）。
\setcounter{enumi}{2}
\item
  对 \(p \geqslant 7\)，证明 \(\mathbf{S}^{p-3}(\mathrm{V})\) 的维数不整除 G 的阶。d) 证明每个单 F{[}G{]}-模都同构于诸模 \(\mathbf{S}^{d}(\mathrm{V})\)（\(0 \leqslant d < p\)）之一（利用习题 27, f)）。
\end{enumerate}

\appendix
\renewcommand{\thechapter}{\arabic{chapter}}
\chapter{无单位元的代数}

在本附录中，k 是一个交换环；诸 k-代数都假设是结合的，但不一定有单位元。

\section{正则理想}

设 A 是一个 k-代数。回忆（III, §1, No.~2, 第 430 页），A 的左理想是 A 的这样一个 k-子模 \(\mathfrak{a}\)：关系 \(a \in A\) 与 \(x \in \mathfrak{a}\) 蕴含 \(ax \in \mathfrak{a}\)。我们同样定义右理想与双边理想的概念。若 A 是一个 k-代数而 \(\mathfrak{a}\) 是 A 的一个双边理想，则在过渡到商时，A 中的乘法在 k-模 A/\(\mathfrak{a}\) 上定义一个 k-代数结构（出处同上）。

若 A 的一个左理想按包含关系是 A 的真左理想所成集合的极大元素，则称它是极大的。

定义 1. —— 设 A 是一个 k-代数，\(\mathfrak{a}\) 是 A 的一个左理想。A 的一个元素 u，若 au − a 属于 \(\mathfrak{a}\)（对每个 \(a \in A\)），则称为模 \(\mathfrak{a}\) 的右单位元。我们称理想 \(\mathfrak{a}\) 是正则的，如果存在模 \(\mathfrak{a}\) 的右单位元。

类似地，我们称 A 的右理想 \(\mathfrak{b}\) 是正则的，如果 A 有一个模 \(\mathfrak{b}\) 的左单位元，即一个元素 v，使得 \(va - a \in \mathfrak{b}\)（对每个 \(a \in A\)）。当一个理想是双边理想时，需要指明它作为左理想还是作为右理想是正则的。

当代数 A 有单位元时，A 的单位元素是模 \(\mathfrak{a}\) 的右单位元（对 A 的每个左理想 \(\mathfrak{a}\)）；这时，A 的每个左（或右）理想都是正则的。另一方面，若 A 是一个所有元素都是幂零元素的 k-代数，则 A 是 A 的唯一的正则（左或右）理想。

设 \(\mathfrak{a}\) 是 A 的一个正则左理想，\(u\) 是模 \(\mathfrak{a}\) 的一个右单位元。若 \(\mathfrak{b}\) 是 A 的一个包含 \(\mathfrak{a}\) 的左理想，则 \(u\) 是模 \(\mathfrak{a}\) 的一个右单位元，从而 \(\mathfrak{b}\) 是正则的。于是，A 的既极大又正则的左理想是 A 的正则真左理想所成集合的极大元素。此外，A 的一个左理想 \(\mathfrak{b}\)，若它包含 \(\mathfrak{a}\)，则是真的，当且仅当 \(u\) 不属于 \(\mathfrak{b}\)。因此，真左理想 \(\mathfrak{b}\)（A 的包含 \(\mathfrak{a}\) 的真左理想）所成之集按包含关系是一个归纳集。集合论 III, §2, No.~4, 第 154 页的定理 2 应用于这个集合就给出下面的结果。

命题 1. —— A 的每个正则真左（相应地，右）理想都包含于一个正则极大左（相应地，右）理想之中。

命题 2. —— 设 A 是一个交换 k-代数。A 的理想 \(\mathfrak{a}\) 是正则极大理想，当且仅当伪环 A/ \(\mathfrak{a}\)（I, §8, No.~1, 第 98 页）是一个体。

A 的元素 u 是模 \(\mathfrak{a}\) 的（右或左）单位元，当且仅当 u 在 A/\(\mathfrak{a}\) 中的典范像是代数 A/\(\mathfrak{a}\) 的单位元素。因此，理想 \(\mathfrak{a}\) 是正则的，当且仅当代数 A/\(\mathfrak{a}\) 有单位元。假设这些条件成立。

映射 \(\mathfrak{b} \mapsto \mathfrak{b}/\mathfrak{a}\) 是从代数 A 的包含 \(\mathfrak{a}\) 的理想所成之集到代数 A/\(\mathfrak{a}\) 的理想所成之集上的一个双射。由于这个代数有单位元，这些就是环 A/\(\mathfrak{a}\) 的理想。最后，由 I, §9, No.~1, 第 115 页的定理 1，环 A/\(\mathfrak{a}\) 是一个体，当且仅当它非零且其仅有的理想是 0 与它自身。命题得证。

例. —— 1) 设 V 是交换域 K 上的一个无限维向量空间。设 \(\operatorname{End}_{\mathrm{K}}^{\mathrm{f}}(\mathrm{V})\) 是 \(\operatorname{End}_{\mathrm{K}}(\mathrm{V})\) 的由有限秩自同态组成的 K-子代数。设 W 是 V 的一个向量子空间，并设 \(\mathfrak{a}_{W}\) 是 \(\operatorname{End}_{\mathrm{K}}^{\mathrm{f}}(\mathrm{V})\) 中核包含 W 的元素所成之集；它是 \(\operatorname{End}_{\mathrm{K}}^{\mathrm{f}}(\mathrm{V})\) 的一个左理想。\(\operatorname{End}_{\mathrm{K}}^{\mathrm{f}}(\mathrm{V})\) 的元素 u 是模 \(\mathfrak{a}_{W}\) 的右单位元，当且仅当我们有 \(u(x)=x\)（对每个 \(x\in W\)）。这样的元素 u 存在，也就是说 \(\mathfrak{a}_{W}\) 是正则的，当且仅当 W 是有限维的。理想 \(\mathfrak{a}_{W}\) 既极大又正则，当且仅当 W 的维数为 1。

*2) 设 \(\mathrm{T}\) 是一个局部紧空间，并设 \(\mathcal{C}_0(\mathrm{T})\) 是交换 \(\mathbf{C}\)-代数，它由从 \(\mathrm{T}\) 到 \(\mathbf{C}\) 的在无穷远处趋于 0 的连续映射组成（Gen.~Top., X, §4, No.~4, 第 316 页）；它有单位元，当且仅当 \(\mathrm{T}\) 是紧的。设 \(\mathrm{F}\) 是 \(\mathrm{T}\) 的一个闭子集，并设 \(\mathfrak{a}_{\mathrm{F}}\) 是 \(\mathcal{C}_0(\mathrm{T})\) 中在 F 上的限制为零函数的元素所成之集；它是 \(\mathcal{C}_{0}(T)\) 的一个理想。\(\mathcal{C}_{0}(T)\) 的元素 u 是模 \(\mathfrak{a}_{F}\) 的单位元，当且仅当我们有 \(u(t)=1\)（对每个 \(t\in F\)）。这样的元素 u 存在，也就是说 \(\mathfrak{a}_{F}\) 是正则的，当且仅当 F 是紧的。映射 \(t\mapsto \mathfrak{a}_{\{t\}}\) 是从 T 到 \(\mathcal{C}_{0}(T)\) 的正则极大理想所成之集上的一个双射（TS, I, §3, n° 2, 第 32 页，推论 1）。假设 T 不是紧的，并用 \(\mathfrak{a}\) 表示 \(\mathcal{C}_{0}(T)\) 中由具有紧支集的函数组成的子集；则 \(\mathfrak{a}\) 是 \(\mathcal{C}_{0}(T)\) 的一个理想，它不包含于 \(\mathcal{C}_{0}(T)\) 的任何正则理想之中。

\begin{enumerate}
\def\labelenumi{\arabic{enumi})}
\setcounter{enumi}{2}
\tightlist
\item
  设 \(L^{1}(\mathbf{R})\) 是局部紧群 R 的卷积代数。回忆（参见 Int., VIII, §4, No.~5, 第 38 页），\(L^{1}(\mathbf{R})\) 是 R 上关于勒贝格测度可积的函数的类所成的空间。两个函数 f 与 g 的类的乘积是由下述公式定义的函数 \(f * g\) 的类
\end{enumerate}

\[
(f * g) (s) = \int_ {- \infty} ^ {+ \infty} f (t) g (s - t) d t
\]

对几乎所有的 \(s \in \mathbf{R}\) 成立。代数 \(\mathrm{L}^1(\mathbf{R})\) 没有单位元。对任何 \(a\)（\(\mathbf{R}\) 中的元素），用 \(\mathfrak{m}_a\) 表示满足下述条件的元素 \(f\)（\(\mathrm{L}^1(\mathbf{R})\) 的元素）所成之集

\[
\int_ {- \infty} ^ {+ \infty} f (t) e ^ {- i a t} d t = 0.
\]

由 TS, II, §3, n° 2, 第 252 页，定理 1，映射 \(a \mapsto \mathfrak{m}_a\) 是从 \(\mathbf{R}\) 到代数 \(\mathrm{L}^1(\mathbf{R})\) 的正则极大理想所成之集上的一个双射。*

\section{单位元的添加}

设 A 是一个 k-代数。在 III, §1, No.~2, 第 431 页中，我们定义了由 A 通过添加单位元素 e 而得到的有单位元代数 \(\widetilde{A}\)。我们把 A 与 \(\widetilde{A}\) 的一个双边理想等同起来。k-模 \(\widetilde{A}\) 是子模 ke 与 A 的直和。

命题 3. —— a) 设 \(\widetilde{\mathfrak{a}}\) 是 \(\widetilde{\mathrm{A}}\) 的一个左理想，使得 \(\widetilde{\mathrm{A}} = \widetilde{\mathfrak{a}} + \mathrm{A}\)。我们置 \(\mathfrak{a} = \widetilde{\mathfrak{a}} \cap \mathrm{A}\)。存在一个元素 \(u\)（\(\mathrm{A}\) 的元素），使得 \(u - e\) 属于 \(\widetilde{\mathfrak{a}}\)。若 \(u\) 是这样一个元素，则它是 \(\mathrm{A}\) 的模 \(\mathfrak{a}\) 的右单位元，并且我们有 \(\widetilde{\mathfrak{a}} = \mathfrak{a} + k(u - e)\)；特别地，理想 \(\mathfrak{a}\) 是正则的。

\begin{enumerate}
\def\labelenumi{\alph{enumi})}
\setcounter{enumi}{1}
\item
  反之，设 \(\mathfrak{a}\) 是 A 的一个正则左理想，\(u\) 是 A 的模 \(\mathfrak{a}\) 的一个右单位元。置 \(\widetilde{\mathfrak{a}} = \mathfrak{a} + k(u - e)\)。则 \(\widetilde{\mathfrak{a}}\) 是 \(\widetilde{\mathrm{A}}\) 的一个左理想，使得 \(\widetilde{\mathrm{A}} = \widetilde{\mathfrak{a}} + \mathrm{A}\)，并且我们有 \(\mathfrak{a} = \widetilde{\mathfrak{a}} \cap \mathrm{A}\)。
\item
  若 \(k\) 是一个体，则条件 \(\widetilde{\mathrm{A}} = \widetilde{\mathfrak{a}} + \mathrm{A}\) 等价于说 \(\widetilde{\mathfrak{a}}\) 不包含于 \(\mathrm{A}\) 之中。
\end{enumerate}

在 a) 的假设下，元素 \(e\)（\(\widetilde{\mathbf{A}}\) 的元素）可以写成 \((e - u) + u\)，其中 \(u \in \mathbf{A}\) 且 \(u - e \in \widetilde{\mathfrak{a}}\)。设 \(x = \lambda e + a\) 是 \(\widetilde{\mathbf{A}}\) 的一个元素，其中 \(\lambda \in k\) 且 \(a \in \mathbf{A}\)；若 \(x\) 属于 \(\widetilde{\mathfrak{a}}\)，则元素 \(y = a + \lambda u = x + \lambda (u - e)\)（\(\mathbf{A}\) 的元素）属于 \(\mathfrak{a}\)，并且我们有 \(x = y - \lambda (u - e)\)；这就证明了等式 \(\widetilde{\mathfrak{a}} = \mathfrak{a} + k(u - e)\)。此外，对每个 \(a\)（\(\mathbf{A}\) 中的元素），元素 \(au - a = a(u - e)\)（\(\mathbf{A}\) 的元素）属于 \(\mathfrak{a}\)，因此 \(u\) 是模 \(\mathfrak{a}\) 的右单位元。这就证明了断言 a)。

设 \(\mathfrak{a}\) 与 \(u\) 如 b) 中所述。由于 \(\widetilde{\mathrm{A}}\) 是 \(\mathrm{A}\) 与 \(k(u - e)\) 的直和，我们有 \(\widetilde{\mathfrak{a}} + \mathrm{A} = \widetilde{\mathrm{A}}\) 与 \(\widetilde{\mathfrak{a}} \cap \mathrm{A} = \mathfrak{a}\)。\(\widetilde{\mathfrak{a}}\) 的每个元素都形如 \(x + \lambda (u - e)\)，其中 \(x \in \mathfrak{a}\) 且 \(\lambda \in k\)。对每个 \(a \in \mathrm{A}\)，我们有

\[
a (x + \lambda (u - e)) = a x + \lambda (a u - a).
\]

由于 u 是 A 的模 \(\mathfrak{a}\) 的右单位元，和 \(ax + \lambda(au - a)\) 属于 \(\mathfrak{a}\)，从而 \(\widetilde{\mathfrak{a}}\) 是 A 的一个左理想。这就证明了 b)。

最后，若 \(k\) 是一个体，则 \(A\) 是 \(\widetilde{A}\) 中的一个超平面，而 \(\widetilde{A}\) 是 \(\widetilde{\mathfrak{a}}\) 与 \(A\) 之和，当且仅当 \(\widetilde{\mathfrak{a}}\) 不包含于 \(A\) 之中。

推论. —— A 的正则左理想是形如 \(A \cap \widetilde{\mathfrak{a}}\) 的理想，其中 \(\widetilde{\mathfrak{a}}\) 是 \(\widetilde{A}\) 的一个左理想，使得 \(\widetilde{A} = \widetilde{\mathfrak{a}} + A\)。

命题 4. —— a) 设 \(\mathfrak{a}\) 是 A 的一个正则极大左理想。存在唯一的左理想 \(\widetilde{\mathfrak{a}}\)（\(\widetilde{\mathrm{A}}\) 的左理想），使得 \(\widetilde{\mathrm{A}} = \widetilde{\mathfrak{a}} + \mathrm{A}\) 且 \(\mathfrak{a} = \widetilde{\mathfrak{a}} \cap \mathrm{A}\)。这个理想是极大的且不包含 A。

\begin{enumerate}
\def\labelenumi{\alph{enumi})}
\setcounter{enumi}{1}
\tightlist
\item
  映射 \(\widetilde{\mathfrak{a}}\mapsto \widetilde{\mathfrak{a}}\cap \mathrm{A}\) 是从 \(\widetilde{\mathrm{A}}\) 的不包含 A 的极大左理想所成之集到 A 的正则极大左理想所成之集上的一个双射。
\end{enumerate}

设 \(\mathfrak{a}\) 是 A 的一个正则极大左理想。由命题 3, a)，左理想 \(\mathfrak{b}\)（\(\widetilde{\mathrm{A}}\) 的左理想，满足 \(\mathfrak{b} + \mathrm{A} = \widetilde{\mathrm{A}}\) 与 \(\mathfrak{b} \cap \mathrm{A} = \mathfrak{a}\)）就是诸理想 \(\mathfrak{a} + k(u - e)\)，其中 \(u\) 是模 \(\mathfrak{a}\) 的右单位元。因此，为证明 \(\widetilde{\mathfrak{a}}\) 的唯一性，只需证明 A 的两个右单位元 \(u\) 与 \(u'\)（模 \(\mathfrak{a}\) 的右单位元）模 \(\mathfrak{a}\) 同余即可。用反证法推理，假设 \(u - u'\) 不属于 \(\mathfrak{a}\)。公式 \(x(u - u') = (xu - x) - (xu' - x)\) 表明我们有 \(\mathrm{A}(u - u') \subset \mathfrak{a}\)；由此推出 \(\mathfrak{a} + k(u - u')\) 是 A 的一个包含 \(\mathfrak{a}\) 且异于 \(\mathfrak{a}\) 的左理想。由于 \(\mathfrak{a}\) 是极大的，我们因此有 \(\mathfrak{a} + k(u - u') = \mathrm{A}\)，从而 \(\mathrm{AA} \subset \mathfrak{a}\)。对每个 \(x \in \mathrm{A}\)，我们有 \(x \equiv xu (\mathrm{mod} \mathfrak{a})\)，因此由上所述 \(x \in \mathfrak{a}\)，这与假设 \(\mathfrak{a} \neq \mathrm{A}\) 矛盾。

于是存在唯一的左理想 \(\widetilde{\mathfrak{a}}\)（\(\widetilde{\mathrm{A}}\) 的左理想），使得 \(\widetilde{\mathrm{A}} = \widetilde{\mathfrak{a}} + \mathrm{A}\) 且 \(\mathfrak{a} = \widetilde{\mathfrak{a}} \cap \mathrm{A}\)。设 \(\mathfrak{b}\) 是 \(\widetilde{\mathrm{A}}\) 的一个包含 \(\widetilde{\mathfrak{a}}\) 的真左理想。则 \(\mathfrak{b} \cap \mathrm{A}\) 是 \(\mathrm{A}\) 的一个包含 \(\mathfrak{a}\) 的真左理想。由于 \(\mathfrak{a}\) 是极大的，它因此等于 \(\mathfrak{a}\)，从而由 VIII, 第 4 页的引理 2 推出 \(\mathfrak{b} = \widetilde{\mathfrak{a}}\)。这就证明了 \(\widetilde{\mathfrak{a}}\) 是 \(\widetilde{\mathrm{A}}\) 的一个极大理想；对这样一个理想，条件 \(\widetilde{\mathrm{A}} = \widetilde{\mathfrak{a}} + \mathrm{A}\) 意味着 \(\widetilde{\mathfrak{a}}\) 不包含 A。

剩下的要证明：若 \(\widetilde{\mathfrak{a}}\) 是 \(\tilde{A}\) 的一个不包含 A 的极大左理想，则 A 的左理想 \(\mathfrak{a} = A \cap \widetilde{\mathfrak{a}}\) 是极大的。设 \(\mathfrak{b}\) 是 A 的一个包含 \(\mathfrak{a}\) 的真左理想。设 u 是模 \(\mathfrak{a}\) 的一个右单位元，使得 \(\widetilde{\mathfrak{a}} = \mathfrak{a} + k(u - e)\)（命题 3）。它也是模 \(\mathfrak{b}\) 的右单位元；用 \(\widetilde{\mathfrak{b}}\) 表示理想 \(\mathfrak{b} + k(u - e)\)（\(\tilde{A}\) 的理想）。由命题 3, b)，我们有 \(\mathfrak{b} = A \cap \widetilde{\mathfrak{b}}\)。理想 \(\widetilde{\mathfrak{a}}\)（等于 \(\mathfrak{a} + k(u - e)\) 的理想）包含于 \(\widetilde{\mathfrak{b}}\)，从而等于 \(\widetilde{\mathfrak{b}}\)，因为 \(\widetilde{\mathfrak{b}}\) 异于 \(\tilde{A}\) 而 \(\widetilde{\mathfrak{a}}\) 极大。从而我们有 \(\mathfrak{a} = \mathfrak{b}\)，故 \(\mathfrak{a}\) 是极大的。

我们把命题 3 与命题 4 向右理想的翻译留给读者。

\section{代数的根}

设 A 是一个 k-代数。A 上的左伪模是一个 k-模 M，它带有伪环 A（II, 附录, No.~2, 第 378 页）上的左伪模结构，使得我们有 \(a(\lambda x) = \lambda(ax) = (\lambda a)x\)（对 \(\lambda \in k\)、\(a \in A\) 与 \(x \in M\)）。我们同样定义 A 上的右伪模。

设 M 是 A 上的一个左伪模；我们用下述方式在 M 上定义一个称为典范的左 \(\tilde{A}\)-模结构

\[
(\lambda e + a) x = \lambda x + a x \quad \text { for } \lambda \in k, a \in A, x \in M.
\]

反之，每个左 \(\widetilde{A}\)-模通过把算子环限制到 k 与 A，都典范地带有 A 上的一个左伪模结构。于是，A 上的左伪模的结构种与 \(\widetilde{A}\) 上的模的结构种是等价的（集合论, IV, §1, No.~7, 第 262 页）。

设 M 是 A 上的一个左伪模，并设 N 是 M 的一个 k-子模。则 N 是 M 的 \(\widetilde{A}\)-子模，当且仅当它在 A 的作用下稳定；这时我们称 N 是 M 的一个子伪模。

如同环的情形，我们定义 A 上的左伪模 \(A_{s}\) 与右伪模 \(A_{d}\)。A 的左（相应地，右）理想是 \(A_{s}\)（相应地，\(A_{d}\)）的子伪模。

设 \(\mathfrak{a}\) 是 A 的一个正则左理想，\(u\) 是模 \(\mathfrak{a}\) 的一个右单位元。置 \(\mathrm{M} = \mathrm{A}_s / \mathfrak{a}\)，并把 \(u\) 在 M 中的像记为 \(z\)。我们有 \(\mathrm{M} = \mathrm{A}z\)，且 \(\mathfrak{a}\) 是 \(z\) 的零化子。

反之，设 M 是 A 上的一个左伪模，并设 z 是 M 的一个元素，使得 M = Az。则存在 A 的一个元素 u，使得 z = uz。对每个 \(a \in A\)，我们有 \((au - a)z = 0\)，故 au − a 属于 z 的零化子 \(\mathfrak{a}\)。于是，\(\mathfrak{a}\) 是 A 的一个正则左理想，元素 u 是模 \(\mathfrak{a}\) 的右单位元，并且在过渡到商时，映射 \(a \mapsto az\) 定义一个从 \(A_{s}/\mathfrak{a}\) 到 M 的同构。

定义 2. —— 我们称 A 上的伪模 M 是单的，如果我们有 AM ≠ 0，并且 0 与 M 是 M 的仅有的一些子伪模。

这等价于说 \(\tilde{A}\)-模 M 是单的且其零化子不包含 A。当 A 是一个环而 M 是一个 A-模时，我们有 AM = M，故定义 2 与 VIII, 第 45 页的定义 1 一致。

命题 5. —— a) 设 \(\mathfrak{m}\) 是 A 的一个正则极大左理想。则伪模 \(A_s / \mathfrak{m}\) 是单的，并且存在 \(A_s / \mathfrak{m}\) 的一个非零元素，其零化子为 \(\mathfrak{m}\)。

\begin{enumerate}
\def\labelenumi{\alph{enumi})}
\setcounter{enumi}{1}
\item
  设 M 是一个单伪模，并设 x 是 M 的一个非零元素。我们有 M = Ax，x 的零化子 \(\mathfrak{m}\) 是 A 的一个正则极大左理想，并且在过渡到商时，映射 \(a \mapsto ax\) 定义一个从 \(A_{s}/\mathfrak{m}\) 到 M 的同构。
\item
  设 M 是一个不约化为 0 的伪模。则 M 是单的，当且仅当对 M 的每个非零元素 x 我们有 M = Ax。
\end{enumerate}

设 \(\mathfrak{m}\) 是 A 的一个正则极大左理想。我们已经看到，存在一个非零元素 \(z\)（\(\mathrm{A}_s / \mathfrak{m}\) 的元素），其零化子等于 \(\mathfrak{m}\) 且使得 \(\mathrm{A}z = \mathrm{A}_s / \mathfrak{m}\)；特别地，我们有 \(\mathrm{A}(\mathrm{A}_s / \mathfrak{m}) \neq 0\)。此外，\(\mathrm{A}_s / \mathfrak{m}\) 的每个子伪模都形如 \(\mathfrak{n} / \mathfrak{m}\)，其中 \(\mathfrak{n}\) 是 A 的一个包含 \(\mathfrak{m}\) 的左理想。由于 \(\mathfrak{m}\) 是极大的，仅有的可能是 \(\mathfrak{n} = \mathfrak{m}\) 与 \(\mathfrak{n} = \mathrm{A}\)，从而伪模 \(\mathrm{A}_s / \mathfrak{a}\) 是单的。这就证明了 a)。

在 b) 的假设下，M 中使得 Ay = 0 的元素 y 所成之集是 M 的一个异于 M 的子伪模，从而约化为 0。于是，Ax 是 M 的一个非零子伪模，由此推出 M = Ax。断言 b) 由定义 2 前的注记得出。

设 M 是一个非零伪模。假设对 M 的每个非零元素 x 我们有 Ax = M。特别地，我们有 AM ≠ 0。设 N 是 M 的一个非零子伪模，并设 x 是 N 的一个非零元素。我们有 Ax = M，因此 N = M。故 M 是单的。反之，若 M 是单的，则由 b)，对每个 \(x \neq 0\) 我们有 M = Ax。

定义 3. —— k-代数 A 的根，记作 \(\Re(\mathrm{A})\)，是 A 的正则极大左理想的交。

当 A 是一个环时，A 的每个左理想都是正则的，故根的定义与 VIII, 第 154 页的定义 2 一致。

例. —— 1) \(k\)-代数 \(\operatorname{End}_k^{\mathrm{f}}(\mathrm{V})\) 的根约化为 0（VIII, 第 436 页，例 1）。*代数 \(\mathcal{C}_0(\mathrm{T})\) 与 \(\mathrm{L}^1 (\mathbf{R})\)) 的根亦然。*

\begin{enumerate}
\def\labelenumi{\arabic{enumi})}
\setcounter{enumi}{1}
\tightlist
\item
  设 A 是一个所有元素都是幂零元素的伪环。A 的根等于 A，因为 A 是唯一的正则左理想。
\end{enumerate}

命题 5 立即给出下面的结果。

命题 6. —— 代数 A 的根是单伪模的零化子的交。它特别是 A 的一个双边理想。

命题 7. —— A 的根是 \(\widetilde{A}\) 的根在 A 上的迹；它也等于反代数 \(A^{\circ}\) 的根（即 A 的正则极大右理想的交）。若环 k 无根，则 A 的根等于 \(\widetilde{A}\) 的根。

等式 \(\Re (\widetilde{\mathrm{A}})\cap \mathrm{A} = \Re (\mathrm{A})\) 由 VIII, 第 438 页，命题 4, b) 得出。由于 \(\widetilde{\mathrm{A}}\) 与 \(\widetilde{\mathrm{A}}^{\circ}\) 有相同的根（VIII, 第 156 页，推论 1），等式 \(\Re (\mathrm{A}) = \Re (\mathrm{A}^{\circ})\) 得出。若 \(k\) 无根，则 \(\widetilde{\mathrm{A}}\) 的包含 A 的极大左理想的交等于 A。于是 \(\Re (\widetilde{\mathrm{A}})\) 包含于 A，从而等于 \(\Re (\mathrm{A})\)。

注. —— 设 x 与 y 是 A 的元素；我们称 x 是 y 的左拟逆元，或 y 是 x 的右拟逆元，如果在 \(\widetilde{A}\) 中，x − e 是 y − e 的左逆，即如果我们有 \(x + y = xy\)。由命题 7 与雅各布森定理（VIII, 第 156 页，定理 1），A 的根由 A 中这样的元素 x 组成：对 \(\widetilde{A}\) 中的每个 u，ux − e 在 \(\widetilde{A}\) 中左可逆。由于我们有 \((a + \lambda e)x - e = ax + \lambda x - e\)（对 \(a \in A\) 与 \(\lambda \in k\)），A 的根是 A 中这样的元素 x 所成之集：\(ax + \lambda x\)（对每个 \(a \in A\) 与每个 \(\lambda \in k\)）在 A 中有左拟逆元。

\section{稠密性定理}

设 A 是一个 k-代数。若左伪模 M 是一族单左伪模的直和，则称它是半单的。

引理 1. —— 设 M 是一个半单左 A-伪模。设 B 是 \(\widetilde{A}\)-模 M 的双交换子。则 M 的每个 A-子伪模都是 M 的 B-子模。

\(\widetilde{\mathrm{A}}\)-模 M 是半单的。设 N 是 M 的一个 A-子伪模。则 N 是 M 的一个 \(\widetilde{\mathrm{A}}\)-子模。由 VIII, 第 56 页的推论 2，存在 A-模 M 的一个以 N 为像的投影子 \(p\)。由于我们有关系 \(pb = bp\)（对每个 \(b \in \mathrm{B}\)），我们得到 N 是 M 的一个 B-子模。

引理 2. —— 设 M 是一个半单左 A-伪模，并设 x 是 M 的一个元素。存在一个元素 \(a \in A\)，使得 ax = x。

设 N 是左 A-伪模 \(\widetilde{Ax}/Ax\)。它满足 \(AN = \{0\}\)。把 VIII, 第 56 页的推论 3 应用于 \(\widetilde{A}\)-模 M 与 \(\widetilde{Ax}\)，得知 A-伪模 N 是半单的。按定义，N 的每个单子伪模 S 满足 \(AS \neq \{0\}\)。于是，伪模 N 是零，并且我们得到 \(x \in Ax\)。

定理 1（雅各布森稠密性定理）. —— 设 M 是一个半单左 A-伪模。设 b 是 \(\widetilde{\mathrm{A}}\)-模 M 的双交换子的一个元素。设 \(\mathrm{F} = \{x_1, \ldots, x_n\}\) 是 M 的一个有限子集。则存在一个元素 \(a \in \mathrm{A}\)，使得 \(bx_i = ax_i\)（对每个 \(i \in [1, n]\)）。

设 B 是 \(\widetilde{\mathrm{A}}\)-模 M 的双交换子。A-伪模 \(\mathrm{M}^n\) 是半单的。设 \(\boldsymbol{x} = (x_1, \ldots, x_n) \in \mathrm{M}^n\)。由引理 2 推出 \(\boldsymbol{x} \in \mathrm{A}\boldsymbol{x}\)。\(\widetilde{\mathrm{A}}\)-模 \(\mathrm{M}^n\) 的双交换子与 B-模 \(\mathrm{M}^n\) 的位似相一致（VIII, 第 79 页，命题 2）。于是由引理 1，A-子伪模 \(\mathrm{A}\boldsymbol{x}\)（\(\mathrm{M}^n\) 的 A-子伪模）是 \(\mathrm{M}^n\) 的一个 B-子模。因此我们有包含关系 \(\mathrm{B}\boldsymbol{x} \subset \mathrm{A}\boldsymbol{x}\)。所以存在一个 \(a \in \mathrm{A}\)，使得 \(b\boldsymbol{x} = a\boldsymbol{x}\)。结论得证。
